% Modernised reading — fonds Alexandre Grothendieck, shelfmark 6, pages 1-66.
%
% This is an INTERPRETATION, not a transcription. It re-reads the pages in
% today's notation and vocabulary, states the hypotheses the manuscript leaves
% implicit, and drops the critical apparatus so that the mathematics reads
% straight through. Everything the brackets and underlines carried — a doubtful
% reading, a notation he used differently, a missing passage — moves into a
% footnote.
%
% It is held to being mathematically correct as it stands, not merely faithful
% to the page. Where the manuscript is loose or mistaken, the true statement is
% what appears here, and the footnote says what the page has. In any dispute of
% substance the transcription governs: whoever cites Grothendieck must cite the
% batch-NN.fr.tex files.
%
% Scope: ONE document for the whole shelfmark, its four batches (pages 1-66)
% read in one pass before any of it was written. The folder is two things: his
% own corrected typed copy of the letter to Tate of May 1966 (pages 1, 3-31),
% and later manuscript notes (pages 32-65; one of them cites Demazure-Gabriel,
% 1970, which agrees with the catalogue's « à partir de 1970 »). What ties them
% is one question, which the letter asks at page 31 (3.2) and the notes come
% back to at pages 40-44: how much of a p-divisible group, or of the
% cohomology of a variety, is captured by a linear object with a Frobenius.
%
% NOTATION fixed once for the folder.
%   - The letter's machine-underlined letters (Z, O, M, Omega...) are written
%     in today's way: \mathbf{Z}, \mathcal{O}, \mathcal{M}, \Omega.
%   - « cristal » in the letter is a crystal on the INFINITESIMAL site — nilpotent
%     thickenings, no divided powers. It is called « cristal (infinitésimal) »
%     here whenever the difference matters, and « cristal cristallin » or
%     « cristal à puissances divisées » is reserved for Berthelot's notion. This
%     is the single most important convention of the reading: the letter's
%     negative results (2.7, 2.8, 3.1) are true of the first and false of the
%     second.
%   - « p-cristal » = F-cristal, « p^{-1}-cristal » = V-cristal, « bicristal de
%     poids i » kept, « cristal de Dieudonné » = bicristal de poids 1,
%     « bi-isocristal » = F-isocristal. The Tate bicrystal T^i: unit crystal
%     with F = V = p^i, weight 2i.
%   - The p-divisible group the typescript writes with the typed sign « Ø » is
%     written Phi; the universal vector extension is E(A), E(Phi) (the letter's
%     G(A), G(Ø)).
%   - In the Tate-Hodge pages the page's [i] (a Tate twist) is written (i); the
%     page's Q_ell, where context forces ell = p, is written Q_p; C there is the
%     completion of an algebraic closure of the local field K.
%
% Substantive departures from the pages, each footnoted where it occurs:
%   - p. 8: VF = p.id (the p is missing on the page); the Tate bicrystal of
%     weight 2i is the i-th power of the one of weight 2 (page: « poids 1 »);
%   - p. 4: « voisinage de R dans S » read « de S dans R »; p. 5: F_R read F_S;
%     p. 12: « stratifiés relativement à X_K » read « à K », and local freeness
%     asserted only on smooth X_K in char. 0, the case used;
%   - p. 11: the category of crystals is abelian only under hypotheses (S smooth
%     over a perfect field, coherent crystals) — said where the page says it
%     without them;
%   - p. 19: spectral sequence with C_{X^an} for C_{S^an};
%   - p. 20: the duality pairing lands in T^d[-2d] (cohomological convention);
%   - p. 21: the relative Frobenius is X_p -> X_p^{(p/S)}, not « X_p^{(p/S)} -> X »;
%     transposition of F by Poincaré duality gives FV = VF = p^d, not p^{2d};
%   - pp. 21-26: the three « affirmations » cannot hold on all nilpotent
%     thickenings (the letter itself shows it in 2.7); they hold on thickenings
%     with locally nilpotent divided powers — Mazur-Messing, Grothendieck-
%     Messing — and that is what is said; 2.6 b) is false in char. p, as the
%     pencil « faux » says;
%   - p. 31: 3.2 is stated under divided powers on the maximal ideal (e < p-1);
%   - pp. 33-36: the leaves are read in the order 33, 34, 36, 35 (p. 36 ends on
%     « Je conjecture », p. 35 opens on « que cet élément… »);
%   - p. 47: « dim 2 au plus » read « dim 1 au plus »; « hypersingulière » (no
%     points of order p) equals supersingular only in dimension <= 2;
%   - p. 59 is struck through on the page; its statements are reported as
%     withdrawn drafts;
%   - p. 65: the Lemma needs k separably closed (his own ink correction of
%     « alg. clos »).
%
% Announced and never established in the folder: the cohomology theory DR(f)
% of Chapter 2 (with Künneth and duality); the conjecture of 2.8 on ordinary
% abelian varieties; the conjecture of p. 36 (left as « Je conjecture » and
% resumed on p. 35 only by its conclusion); b) of p. 62 (homomorphisms étale ->
% multiplicative type, a heading only); the Witt-vector version of p. 65, which
% breaks off. Page 38 is Tate's, not his, and is only described.
%
% Pass: Opus 5.5 (claude-opus-5-5), 2026-10-03 — first pass, unchecked against the pages by a human.
\documentclass[11pt,a4paper]{article}
\input{../preamble/grothendieck}

\folder{6}
\pages{1}{66}
\dating{1966-[à partir de 1970]}
\watermark{Édition de démonstration\\interprétation personnelle de l'œuvre}
\foldertitle{Lettre Tate (mai 66) (Cristaux) : lettre (1966), tapuscrit (s.d.), notes manuscrites (s.d.) — lecture modernisée du dossier entier}

\begin{document}

\section*{Résumé}

\begin{resume}
Ce dossier s'ouvre sur une lettre. En mai 1966, de Pise, Grothendieck écrit à
John Tate qu'il a « réfléchi aux groupes formels et à la cohomologie de De
Rham » et qu'il est arrivé à « un début de théorie » qu'il a envie d'exposer
« pour se clarifier les idées ». La lettre fait trente pages dactylographiées ;
le dossier en garde une copie qu'il a reprise à l'encre et au crayon. Elle est
suivie de notes manuscrites plus tardives, écrites au fil de plusieurs années,
qui reprennent les mêmes questions avec d'autres outils.

Le problème est le suivant. Pour une variété définie sur les nombres
complexes, la cohomologie — les « trous » de la variété, mesurés par des
nombres et des espaces vectoriels — se calcule de plusieurs façons, et ces
façons s'accordent. En caractéristique $p$, là où l'on compte modulo un nombre
premier, l'outil analytique le plus naturel, la cohomologie de De Rham (celle
des formes différentielles), se comporte mal : elle donne des espaces
vectoriels sur un corps de caractéristique $p$, qui perdent l'information
fine, et elle varie mal en famille. On voudrait une cohomologie à valeurs dans
un anneau de caractéristique nulle — les vecteurs de Witt, qui sont à $\mathbf{F}_p$
ce que les entiers $p$-adiques sont aux entiers modulo $p$ — et qui ne dépende
que de la variété en caractéristique $p$, pas d'un relèvement choisi.

L'idée de la lettre est dans son premier paragraphe, et elle est donnée par
une image. Un cristal, dit-il, a deux propriétés : la \emph{rigidité}, et la
faculté de \emph{croître} dans un milieu approprié. Un objet « cristallin » sur
un espace $S$ est un objet qui, chaque fois qu'on épaissit $S$
infinitésimalement — qu'on lui ajoute des directions nilpotentes, comme on
passe de $\mathbf{F}_p$ à $\mathbf{Z}/p^n$ —, s'étend à l'épaississement, et
de façon unique : il croît sans se déformer. Si la cohomologie d'une variété
en caractéristique $p$ est un tel objet, alors sa valeur sur les
épaississements $p$-adiques est une cohomologie à coefficients entiers
$p$-adiques, indépendante de tout choix. C'est la naissance de ce qu'on
appelle aujourd'hui la cohomologie cristalline. La lettre ajoute une structure
qui fait toute la richesse de la théorie : l'action du Frobenius, l'élévation
à la puissance $p$, qui transforme ces objets linéaires en « modules de
Dieudonné » — les invariants linéaires, connus depuis Dieudonné, des groupes
formels en caractéristique $p$.

La lettre est aussi le récit d'un échec partiel, et c'est une des raisons de la
lire. Au milieu du chapitre 2, Grothendieck écrit : « je m'aperçois avec
consternation » que ses « brillantes conjectures » sont fausses telles
quelles ; il donne le contre-exemple — les courbes elliptiques dites
supersingulières — et conclut que sa notion de cristal est « irrémédiablement
trop fine ». Le diagnostic est juste, et une note au crayon dans la marge en
donne déjà le remède : « faut introduire des puissances divisées ». C'est ce
que fera Pierre Berthelot dans sa thèse : en n'épaississant que d'une façon
qui permette de diviser $x^n$ par $n!$, on obtient une théorie où les
affirmations de la lettre deviennent des théorèmes. On voit ici, à découvert,
comment une théorie naît : une idée juste, une première définition un peu trop
rigide, un contre-exemple honnêtement constaté, une correction. « Il est
probable qu'il y aura bien d'autres canulars de détail dans ces notes, mais,
je pense, sans conséquence ! » — la phrase est de la lettre.

Les notes qui suivent changent d'angle. Elles regardent, pour une courbe
elliptique sur un corps $p$-adique, comment le groupe de Galois agit sur ses
points de torsion d'ordre une puissance de $p$, et posent les questions qui
fondent ce qu'on appelle aujourd'hui la théorie de Hodge $p$-adique ; elles
reprennent une question de Barsotti sur les variétés abéliennes que
distinguent ou non leurs groupes formels ; elles calculent l'espace des
déformations d'une variété abélienne « ordinaire », qui se trouve être un
tore ; elles construisent enfin, pour les groupes finis en caractéristique
$p$, la théorie de Dieudonné dans ses deux cas extrêmes, étale et
multiplicatif. Une question les relie à la lettre, posée à sa dernière page :
comment lire, sur l'objet linéaire muni de son Frobenius et de sa filtration,
l'action de Galois sur les points de torsion ? C'est le problème que
Grothendieck appellera plus tard celui de la « fonction mystérieuse ».

Pour aller plus loin, on cherchera : cohomologie cristalline, site
infinitésimal, $F$-cristaux et $F$-isocristaux, cristal de Dieudonné d'un
groupe $p$-divisible, théorème de Grothendieck--Messing, coordonnées de
Serre--Tate, décomposition de Hodge--Tate, théorie de Hodge $p$-adique.

\keywords{crystal, infinitesimal site, crystalline cohomology, divided powers, stratification, Cartier descent, F-crystal, Dieudonné crystal, F-isocrystal, Tate twist, rigid analytic space, Gauss-Manin connection, Monsky-Washnitzer cohomology, universal vector extension, Grothendieck-Messing theory, Serre-Tate theorem, Manin map, supersingular elliptic curve, canonical lift, p-divisible group, Tate module, Tate curve, Hodge-Tate decomposition, p-adic Hodge theory, Sen theory, Barsotti isogeny question, Serre-Tate coordinates, Witt vectors, Dieudonné module, truncated Barsotti-Tate group, co-Lie complex, unit-root F-crystal, Artin-Schreier sequence, Lang's theorem}
\end{resume}

\pagerange{1}{1}
\section*{Le fil du dossier, et les conventions}

\textbf{Le fil.} Le dossier a deux parties, de dates différentes.

\begin{itemize}
\item \emph{La lettre à Tate, Pise, mai 1966} (pages 1 et 3 à 31). Trois
chapitres. Le chapitre 1 (pages 1 à 18) définit les cristaux, leurs variantes
munies d'un Frobenius, les isocristaux, et les décrit dans les cas
calculables ; il propose un site dont les cristaux sont les faisceaux
quasi-cohérents, et formule le « test-clef » : la cohomologie de ce site
doit donner la cohomologie de De Rham. Le chapitre 2 (pages 18 à 29) postule
une cohomologie de De Rham à valeurs dans les cristaux, l'explicite pour les
schémas abéliens (trois « affirmations »), puis la réfute sous la forme
proposée (2.7) et en tire les conséquences (2.8). Le chapitre 3 (pages 30 et
31) transpose la construction aux groupes $p$-divisibles et s'arrête sur une
question.
\item \emph{Les notes manuscrites} (pages 32 à 65), sous des chemises
titrées de sa main : la réduction des courbes elliptiques et leur module de
Tate $p$-adique (pages 32 à 36) ; une feuille de Tate (pages 37 et 38) ; des
questions de « théorie de Tate--Hodge » (pages 39 à 44) ; la question
d'isogénie de Barsotti (pages 46 à 50) ; les déformations d'extensions de
groupes de Barsotti--Tate et la structure locale des schémas modulaires
(pages 51 à 56) ; des compléments sur les groupes de Witt (pages 58 et 59) ;
les groupes finis en caractéristique $p$ et leurs cristaux de Dieudonné
(pages 60 à 65).
\end{itemize}

La question de la page 31 — déterminer le module de Tate de la fibre générique
d'un groupe $p$-divisible à partir de son module de Dieudonné filtré — est
celle que les pages 40 à 44 reprennent du côté galoisien. Les pages 51 à 65
reprennent, du côté des groupes, le cristal de Dieudonné que la lettre
cherchait à atteindre par la cohomologie de De Rham.

\textbf{Deux notions de cristal.} Le mot « cristal » a aujourd'hui deux sens,
et la lettre n'en connaît qu'un. Ses cristaux sont définis sur les
épaississements \emph{nilpotents} quelconques ; on les appelle aujourd'hui
cristaux sur le \emph{site infinitésimal}. Ceux de la théorie cristalline,
due à Berthelot, ne sont définis que sur les épaississements munis de
\emph{puissances divisées}. Les deux théories coïncident en caractéristique
nulle et divergent en caractéristique $p$. On dit ici « cristal » au sens de
la lettre, et « cristal à puissances divisées » pour l'autre ; là où la
différence décide de la vérité d'un énoncé — c'est le cas aux pages 13, 17,
21 à 31 —, on le dit.\footnote{Le crayon de la page 17 l'indique déjà,
« (faut d'introduire des puiss.\ div.) », de lecture en partie conjecturale.
Le site cristallin, avec puissances divisées, est celui de P.~Berthelot,
\emph{Cohomologie cristalline des schémas de caractéristique $p>0$}, Lecture
Notes in Math.\ 407 (1974), issu de sa thèse de 1972 ; Grothendieck avait
exposé le site infinitésimal et esquissé le site cristallin à l'IHÉS en
décembre 1966 (\emph{Crystals and the de Rham cohomology of schemes}, notes de
J.~Coates et O.~Jussila, dans \emph{Dix exposés sur la cohomologie des
schémas}, 1968).}

\textbf{Noms.} La lettre parle de $p$-cristal, de $p^{-1}$-cristal, de
bicristal de poids $i$ et de bi-isocristal. On dit aujourd'hui
$F$-cristal, $V$-cristal, et $F$-isocristal pour le dernier ; on garde
« bicristal de poids $i$ », qui n'a pas d'équivalent courant, et « cristal de
Dieudonné » pour le poids 1, nom que la lettre propose elle-même. On note
$\Phi$ le groupe $p$-divisible que le tapuscrit écrit avec le signe
« Ø »,\footnote{Le signe tapé « Ø » sert, page 23, à désigner une filtration,
et, pages 30 et 31, un groupe $p$-divisible ; on écrit $\Phi$ pour le second et
on ne donne pas de lettre à la première.} et $E(A)$, $E(\Phi)$ l'extension
vectorielle universelle que la lettre note $G(A)$, $G(\Phi)$.

\textbf{L'ordre des feuillets.} Il est celui du fac-similé, sauf en un point :
les pages 33 à 36 se lisent dans l'ordre 33, 34, 36, 35 (la page 36 s'arrête
sur « Je conjecture », la page 35 commence par « que cet élément a une image
non nulle »). Le crayon des archivistes numérote les pages 39 à 66 du
fac-similé de 50 à 77 ; on suit le fac-similé.

\pagerange{1}{3}
\section*{La lettre, chapitre 1 : la notion de cristal (pages 1 à 18)}

\subsection*{Les épaississements de $S$ et la définition (pages 1 à 3)}

La lettre est une copie dactylographiée, corrigée de sa main à l'encre et au
crayon, datée en tête « Pise, Mai 1966 ».\footnote{Le papier est à en-tête de
l'Istituto Matematico « Leonida Tonelli » de Pise ; la page 2 en est le verso
blanc. La lettre a sa propre pagination, en retard d'une unité sur celle des
archivistes à partir de la page 3. Les corrections à l'encre sont de sa main
(vérification par sondage sur la page 21, Michel Hua, 27 septembre 2026), à une
exception près, vérifiée sur la page 1 : le numéro « 1.0. » de la marge n'est
pas de lui. On ne signale ci-dessous les corrections manuscrites que lorsque le
sens en dépend.} Elle commence par un commentaire terminologique : un cristal
possède deux propriétés caractéristiques, la \emph{rigidité} et la faculté de
\emph{croître} dans un voisinage approprié, et il y a des cristaux de toute
espèce de substance — « de soude, de soufre, de modules, d'anneaux, de schémas
relatifs ».

\textbf{1.0.} Soit $S$ un schéma au-dessus d'un schéma $R$ (le cas principal
est $R = \operatorname{Spec}\mathbf{Z}$). On considère la catégorie $C$ des
\emph{épaississements nilpotents} de $S$ au-dessus de $R$ : les $R$-schémas
$T$ munis d'un $R$-morphisme $S \to T$ qui est une immersion fermée définie par
un idéal localement nilpotent ; les morphismes sont les $R$-morphismes sous
$S$.\footnote{C'est la première version de la définition ; la lettre la
corrige en 1.16 (page 16) en remplaçant $S$ par ses ouverts. On donne ici la
version corrigée dès le paragraphe consacré au site (pages 16 et 17), comme la
lettre le demande.} La catégorie fibrée des modules quasi-cohérents sur des
schémas variables induit sur $C$ une catégorie fibrée, qui associe à $T$ la
catégorie des $\mathcal{O}_T$-modules quasi-cohérents.

\textbf{Définition 1.1.} Un \emph{cristal de modules} (quasi-cohérents) sur
$S$ relativement à $R$ est une section cartésienne de cette catégorie fibrée
au-dessus de $C$ : la donnée, pour tout épaississement $T$, d'un module
quasi-cohérent $\mathcal{M}_T$, et pour toute flèche $u : T \to T'$ de $C$ d'un
\emph{isomorphisme} $u^*\mathcal{M}_{T'} \simeq \mathcal{M}_T$, ces isomorphismes
étant compatibles à la composition. Plus généralement, pour toute catégorie
fibrée $\mathcal{F}$ sur $(\mathrm{Sch})_{/R}$, on définit de même les
$\mathcal{F}$-cristaux : cristaux d'algèbres, d'algèbres commutatives, de
schémas relatifs. Pour $R = \operatorname{Spec}\mathbf{Z}$ on parle de
cristaux \emph{absolus}.

Le mot dit l'idée : un cristal est un objet sur $S$ qui « croît » le long de
tout épaississement, de façon unique à isomorphisme unique près, et de façon
« rigide », puisque toutes les flèches de transition sont des isomorphismes.
Ce sont, en termes d'aujourd'hui, les cristaux quasi-cohérents sur le
\emph{site infinitésimal} de $S$ relativement à $R$.

\textbf{1.2 et 1.3.} Les $\mathcal{F}$-cristaux forment une catégorie qui
dépend fonctoriellement de $\mathcal{F}$, ce qui munit les cristaux de modules
des opérations tensorielles habituelles. Si $S \to R' \to R$, un cristal
relatif à $R$ en définit un relatif à $R'$ par restriction ; en particulier un
cristal absolu définit un cristal relatif pour toute base. Pour $S' \to S$, on
définit l'image inverse en construisant un foncteur $C' \to C$ : à un
épaississement $S' \to T'$ on associe la somme amalgamée
$T = S \amalg_{S'} T'$, qui est un épaississement de $S$. Les cristaux de type
$\mathcal{F}$ forment ainsi une catégorie fibrée sur $(\mathrm{Sch})_{/R}$,
et les deux variances, en $R$ et en $S$, se combinent en une variance en
$(R,S)$.\footnote{La lettre reconnaît en 1.16 que cette construction n'est
raisonnable que pour $S$ et $S'$ affines : la somme amalgamée de schémas
peut ne pas exister en général, et il faut se localiser. Pour les modules
quasi-cohérents, qui forment un champ pour la topologie de Zariski, la
localisation ne pose pas de problème.}

\pagerange{3}{5}
\subsection*{La valeur en $S$ et trois exemples (pages 3 à 5)}

\textbf{1.4.} On a le foncteur « valeur en $S$ », $\mathcal{M} \mapsto
\mathcal{M}_S$, des cristaux de modules vers les $\mathcal{O}_S$-modules
quasi-cohérents. Il n'est en général \emph{pas fidèle} (exemple 1 ci-dessous) :
il est donc dangereux de voir un cristal comme un module sur $S$ muni d'une
structure supplémentaire. Ce point de vue devient raisonnable lorsque le
foncteur est fidèle, ce qui arrive dans les cas de 1.8.

\textbf{Exemple 1 (1.5) : vecteurs de Witt.} Soit $S = \operatorname{Spec}k$,
$k$ un corps parfait, et $R = \operatorname{Spec}\mathbf{Z}$. Soit $W$ l'anneau
de Cohen de $k$ : si $\operatorname{car}k = p > 0$, c'est l'anneau $W(k)$ des
vecteurs de Witt ; si $\operatorname{car}k = 0$, c'est $k$ lui-même, et on
suppose alors $k$ \emph{algébrique} sur $\mathbf{Q}$.\footnote{Cette hypothèse
est ajoutée à l'encre, de sa main, et elle est nécessaire : pour un corps
$k$ de caractéristique nulle transcendant sur $\mathbf{Q}$, un épaississement
local artinien de $k$ admet plusieurs corps de représentants, et $k$ n'est
plus objet final de $C$.} La catégorie $C$ est alors équivalente à celle des
$W$-algèbres locales, annulées par une puissance de l'idéal maximal de $W$, à
corps résiduel $k$.\footnote{Pour $k$ parfait de caractéristique $p$, une telle
algèbre $A$ reçoit un unique homomorphisme $W(k) \to A$ relevant
l'identité de $k$ ; c'est la propriété universelle des vecteurs de Witt sur un
corps parfait, et $p$ est nilpotent dans $A$ puisqu'il est dans l'idéal
maximal, qui est nil.} Écrivant $W = \varprojlim W_n$, on trouve, pour
$p > 0$, qu'un cristal de modules absolu sur $k$ équivaut à un système
projectif $(M_n)$ de $W_n$-modules avec $M_{n+1} \otimes W_n \simeq M_n$, et
que les cristaux de présentation finie forment une catégorie équivalente à
celle des \emph{$W$-modules de type fini}. Pour $p = 0$ on trouve les
$k$-espaces vectoriels. La valeur en $S$ s'identifie alors au foncteur
$M \mapsto M \otimes_W k$, qui, pour $p > 0$, n'est pas fidèle : la
multiplication par $p$ sur $W$ devient nulle.

C'est le premier exemple de ce que la notion veut capter : un objet sur le
corps $k$ de caractéristique $p$ dont la vraie nature est un module sur
l'anneau $p$-adique $W(k)$.

\textbf{Exemple 2 (1.6) : $S$ étale sur $R$.} Si $S = R$, ou plus généralement
si $S$ est étale sur $R$, alors $S$ est objet final de $C$, et la valeur en
$S$ est une équivalence : un cristal est simplement un objet sur $S$. En
particulier, un cristal de modules sur $\operatorname{Spec}\mathbf{Z}$ est un
$\mathbf{Z}$-module.

\textbf{Exemple 3 (1.7) : $S$ sous-schéma de $R$.} On peut supposer $S$ fermé
dans $R$, défini par un idéal quasi-cohérent $J$. Soit $S_n = V(J^{n+1})$ le
$n$-ième voisinage infinitésimal de $S$ dans $R$.\footnote{La page écrit
« voisinage infinitésimal de $R$ dans $S$ » : coquille.} La famille des $S_n$
est finale dans $C$, localement sur $S$,\footnote{Précision ajoutée : un
épaississement $T$ ne se factorise par un $S_n$ que localement, là où l'idéal
de $S$ dans $T$ est nilpotent d'un ordre fixé et $J$ de type fini ; c'est le
cas si $R$ est localement noethérien.} de sorte qu'un cristal équivaut à une
suite compatible d'objets sur les $S_n$. Si $R$ est localement noethérien, un
cristal de modules de présentation finie s'identifie à un module cohérent sur
le complété formel de $R$ le long de $S$. Grâce à Cohen, un cristal absolu
sur $k$ parfait est la même chose qu'un cristal relatif à $W$ : cet exemple
contient donc l'exemple 1 en caractéristique $p$. Les exemples 2 et 3 ont une
généralisation commune, pour $S \to R$ non ramifié, qui permet encore de
construire des voisinages infinitésimaux.

\pagerange{5}{6}
\subsection*{Cristaux et stratifications (pages 5 et 6)}

\textbf{1.8.} Pour $n \geq 0$, soit $\Delta_n$ le $n$-ième voisinage
infinitésimal de la diagonale de $S \times_R S$, avec ses deux projections
$\mathrm{pr}_1, \mathrm{pr}_2$ vers $S$. Pour un objet $E$ de
$\mathcal{F}_S$ :
\begin{itemize}
\item une \emph{$n$-connexion} sur $E$ est un isomorphisme
$\mathrm{pr}_1^*E \simeq \mathrm{pr}_2^*E$ sur $\Delta_n$ induisant l'identité
sur la diagonale ;
\item une \emph{pseudo-stratification} (ou $\infty$-connexion) est une famille
compatible de $n$-connexions pour tout $n$ ;
\item une \emph{stratification} est une pseudo-stratification qui vérifie la
condition de cocycle d'une donnée de descente sur les voisinages
infinitésimaux de la diagonale de $S \times_R S \times_R S$.
\end{itemize}
Si $S$ est formellement non ramifié sur $R$, c'est-à-dire si
$\Omega^1_{S/R} = 0$, ces notions sont triviales : tout objet porte une unique
stratification, et les objets stratifiés s'identifient, par la valeur en $S$, à
$\mathcal{F}_S$ lui-même.\footnote{La page écrit « à la catégorie
$\mathcal{F}_R$ elle-même » ; il faut lire $\mathcal{F}_S$.} Le foncteur
d'oubli des objets stratifiés vers $\mathcal{F}_S$ est toujours fidèle.

Tout cristal $\mathcal{M}$ définit sur sa valeur $\mathcal{M}_S$ une
\emph{stratification canonique} — les deux images inverses de $\mathcal{M}$ sur
$\Delta_n$ sont toutes deux la valeur de $\mathcal{M}$ sur l'épaississement
$\Delta_n$ de $S$ —, d'où un foncteur des cristaux vers les objets stratifiés,
qui d'après 1.4 n'est pas toujours fidèle. C'est une \emph{équivalence} lorsque
$S \to R$ est formellement lisse, par exemple lisse, ou lorsque $S$ et $R$ sont
les spectres de corps $k \supset k_0$ avec $k$ séparable sur $k_0$ ; si $S \to R$
est formellement étale, les deux catégories sont équivalentes à
$\mathcal{F}_S$, ce qui redonne l'exemple 1 en caractéristique nulle.

La lettre avance, « sauf erreur » et sans démonstration, que l'équivalence vaut
encore pour $S \to R$ plat et localement de présentation finie, du moins pour
$R$ localement noethérien et les modules cohérents. On ne reprend pas cet
énoncé comme établi.\footnote{La réserve est de la lettre (« je n'ai pas écrit
de démonstration ») et l'ajout à l'encre la renforce. En général, les
cristaux sur le site infinitésimal de $S$ s'identifient aux modules munis
d'une stratification sur le complété d'un $R$-schéma \emph{lisse} dans lequel
$S$ se plonge, non sur $S$ lui-même ; c'est ce qui rend l'énoncé délicat.}

\pagerange{6}{7}
\subsection*{Digression : la caractéristique nulle et le théorème de Cartier (page 6)}

\textbf{1.9.} Si $S$ est lisse sur $R$ et $R$ de caractéristique nulle, une
stratification sur un module $M$ est déterminée par la $1$-connexion qu'elle
définit, c'est-à-dire par une connexion $\nabla$, et une connexion provient
d'une stratification si et seulement si sa \emph{courbure}, section de
$\Omega^2_{S/R} \otimes \underline{\mathrm{End}}(M)$, est nulle ; autrement
dit, stratifications et connexions intégrables coïncident.\footnote{La lettre
ajoute qu'il suffit que $S$ soit « différentiellement lisse », la diagonale
étant une immersion régulière, et se demande si la lissité est nécessaire et si
l'énoncé s'étend à toute catégorie fibrée.} L'hypothèse de caractéristique
nulle est essentielle. En caractéristique $p > 0$, l'énoncé qui la remplace est
le théorème de Cartier : la donnée d'une connexion intégrable \emph{de
$p$-courbure nulle} sur $M$ équivaut à une donnée de descente de $M$ le long
du Frobenius relatif $S \to S^{(p/R)}$.\footnote{La lettre dit « connexion sans
torsion », c'est-à-dire de courbure nulle, puis ajoute à l'encre « compatible
avec puissances $p$-ièmes » : c'est la nullité de la $p$-courbure, au sens
qu'on donne aujourd'hui à ce mot (le terme est postérieur ; la notion est dans
Cartier).} « Il y a loin de là à une stratification ! » : en caractéristique
$p$, une stratification demande l'action de tous les opérateurs
$\partial^n/n!$, et non des seules dérivations.

\pagerange{6}{9}
\subsection*{$F$-cristaux, $V$-cristaux, cristaux de Dieudonné (pages 6 à 9)}

\textbf{1.10.} On suppose maintenant $S$ de caractéristique $p > 0$ et
$R = \operatorname{Spec}\mathbf{Z}$ (ou $\operatorname{Spec}\mathbf{Z}_p$, ce
qui revient au même). Si $\mathcal{M}$ est un cristal de modules de
présentation finie, sa valeur en un point parfait $s'$ au-dessus d'un point
$s$ de $S$ est, par 1.5, un module de type fini sur $W(k(s'))$ : un cristal est
une « famille algébrique » de modules sur les anneaux de Witt des points
parfaits de $S$ — et quelque chose de plus fin lorsque $S$ n'est pas parfait,
par exemple le spectre d'un corps de fonctions.

Le Frobenius absolu $\mathrm{frob}_S : S \to S$ (puissance $p$-ième) permet
d'associer à tout cristal $\mathcal{M}$ son image inverse
$\mathcal{M}^{(p)} = \mathrm{frob}_S^*\mathcal{M}$. On appelle
\begin{itemize}
\item \emph{$F$-cristal} (la lettre : $p$-cristal) un cristal muni d'un
morphisme $F : \mathcal{M}^{(p)} \to \mathcal{M}$ ;
\item \emph{$V$-cristal} (la lettre : $p^{-1}$-cristal) un cristal muni d'un
morphisme $V : \mathcal{M} \to \mathcal{M}^{(p)}$ ;
\item \emph{bicristal de poids $i$}, lorsque $\mathcal{F}$ est fibrée en
catégories additives, un cristal muni de $F$ et de $V$ avec
\[
  FV = p^i\,\mathrm{id}_{\mathcal{M}}, \qquad VF = p^i\,\mathrm{id}_{\mathcal{M}^{(p)}} .
\]
\end{itemize}
Pour $i = 1$, c'est un \emph{cristal de Dieudonné}.\footnote{La page donne, pour
le poids 1, les relations $FV = p\,\mathrm{id}_M$ et $VF = \mathrm{id}_{M^{(p)}}$ ;
le facteur $p$ manque dans la seconde, et la formule générale qui suit, comme
1.12, exige $VF = p$. La précision « de poids 1 » est ajoutée dans
l'interligne.} Les $F$-cristaux de modules ont les opérations tensorielles
covariantes ; un foncteur contravariant (le dual) échange $F$-cristaux et
$V$-cristaux, et les bicristaux de poids $i$ sont stables par dualité (pour des
objets localement libres de type fini) : c'est pour avoir une notion
« autoduale » que la lettre introduit le bicristal. Le poids 1 lui semble
adapté aux groupes formels, ou, « ce qui revient moralement au même », à la
cohomologie de De Rham en degré 1 ; le poids $i$ au degré $i$.

Le \emph{bicristal de Tate} $T^i$ est le cristal unité $T \mapsto \mathcal{O}_T$
muni de $F = V = p^i$ ; il est de poids $2i$, et $T^i$ est la puissance
tensorielle $i$-ième de $T^1$, qui est de poids 2.\footnote{La page dit « la
puissance tensorielle $i$-ème du bicristal de Tate de poids 1 » ; c'est le
poids 2 qu'il faut lire, celui de $T^1$ (les poids s'ajoutent par produit
tensoriel, comme la lettre le note aussitôt). Le « $2i$ » est repris à l'encre
dans la marge.} Les bicristaux de poids quelconques se multiplient
tensoriellement, les poids s'additionnant. En termes modernes, $T^1$ est le
$F$-cristal de Tate, celui de la cohomologie $H^2$ de la droite projective.

\pagerange{8}{11}
\subsection*{Isocristaux et $F$-isocristaux (pages 8 à 11)}

\textbf{1.11.} Pour rendre $T^1$ inversible, il faut passer à une catégorie de
fractions. La catégorie des cristaux de modules sur $S$ est
$\mathbf{Z}_p$-linéaire (les $\mathrm{Hom}$ sont des limites projectives de
modules annulés par des puissances de $p$) ; rendre ses $\mathrm{Hom}$
$\mathbf{Q}$-linéaires ou $\mathbf{Q}_p$-linéaires revient au même. Pour les
cristaux de type fini, c'est garder les mêmes objets et remplacer
$\mathrm{Hom}$ par $\mathrm{Hom} \otimes_{\mathbf{Z}} \mathbf{Q}$ ; dans les
cas où la catégorie des cristaux est abélienne, cela revient à quotienter par
la sous-catégorie épaisse des objets de torsion.\footnote{La page dit
« notre catégorie abélienne » sans hypothèse. Faute de platitude des
changements de base, la catégorie des cristaux quasi-cohérents n'a pas de
noyaux en général ; elle est abélienne pour les cristaux cohérents lorsque $S$
est lisse sur un corps parfait, par la description de 1.13, qui est le cas où
l'énoncé sert. La description par $\mathrm{Hom} \otimes \mathbf{Q}$ est
l'ajout à l'encre au pied de la page 8.} On obtient les \emph{isocristaux}
(« cristaux de modules à isogénie près »)\footnote{La page écrit
« isogristaux » : coquille.}, puis les $F$-isocristaux, munis de
$F : \mathcal{M}^{(p)} \to \mathcal{M}$.

Il n'y a plus lieu de parler de bi-isocristal de poids $i$, puisque $V$ serait
déterminé par $F$ comme $p^iF^{-1}$. La lettre appelle donc
\emph{bi-isocristal} un $F$-isocristal dont le $F$ est un \emph{isomorphisme},
sans choisir de poids : c'est ce qu'on appelle aujourd'hui un
\emph{$F$-isocristal}. La notion est autoduale sur les objets de type fini :
le dual de $(\mathcal{M}, F)$ est $(\mathcal{M}^\vee, {}^tF^{-1})$.\footnote{La
page écrit $(\underline{M}, {}^tF^{-1})$ sans marquer le dual, puis « où
$\underline{M}$ est l'isocristal dual » : le signe du dual manque. Une ligne
biffée commençait par « si $\phi$ est un foncteur contravariant ».}

En localisant les bicristaux de poids $i$, on obtient une catégorie
$\mathrm{IsBicr}(S,i)$ et un foncteur exact « oubli de $V$ »
\[
  \mathrm{IsBicr}(S,i) \longrightarrow \mathrm{IsBicr}(S)
\]
vers les $F$-isocristaux. Il est \emph{pleinement fidèle}, puisque $V$ y est
déterminé par $F$. Les foncteurs
$\mathrm{Bicr}(S,i) \to \mathrm{Bicr}(S,i+1)$,
$(\mathcal{M}, F, V) \mapsto (\mathcal{M}, F, pV)$, deviennent après
localisation des foncteurs pleinement fidèles
$\mathrm{IsBicr}(S,i) \to \mathrm{IsBicr}(S,i+1)$ : les bicristaux de poids
croissant forment une généralisation de plus en plus large des cristaux de
Dieudonné, et les $F$-isocristaux les englobent tous, avec « la louable vertu »
d'être stables par produit tensoriel et par dualité. C'est la catégorie que la
lettre propose comme catégorie de valeurs d'une cohomologie de De Rham « à
isogénie près ».

\textbf{Exemple 1 (1.12) : $S$ spectre d'un corps parfait $k$.} Un cristal de
type fini est un $W$-module de type fini $M$ ; le Frobenius $\sigma = f_W$ de
$W$ donne $M^{(p)} = M \otimes_{W,\sigma} W$. Une structure de $F$-cristal
est un endomorphisme $\sigma$-semi-linéaire $F : M \to M$ ; une structure de
$V$-cristal, un endomorphisme $\sigma^{-1}$-semi-linéaire $V$ ; et $(F,V)$
définit un bicristal de poids $i$ si et seulement si $FV = VF = p^i$. Pour
$i = 1$ on retrouve les \emph{modules de Dieudonné} : les cristaux de
Dieudonné sur $k$ sont les modules de Dieudonné. Les isocristaux de type fini
sont les espaces vectoriels de dimension finie sur $K = \operatorname{Frac}W$,
et un $F$-isocristal est un tel espace muni d'une bijection
$\sigma$-semi-linéaire $F$.

L'image essentielle de $\mathrm{IsBicr}(k,i)$ est formée des $(E, F)$ tels
que, posant $V = p^iF^{-1}$, l'orbite de tout vecteur sous les monômes non
commutatifs en $F$ et $V$ soit bornée — « pure tautologie », car un réseau
stable existe alors. Il existe un $i$ convenable si et seulement si les
$F^nx$ sont bornés pour tout $x$.\footnote{Justification, qui n'est pas sur la
page : un monôme en $F$ et $V$ s'écrit $p^{ib}F^{a-b}$ ; si les $F^nx$ sont
bornés, il reste à borner les $(p^iF^{-1})^n x$, ce qui a lieu pour $i$ assez
grand en dimension finie. Dans le langage des pentes de Dieudonné--Manin
(Manin, 1963), la condition dit que les pentes de $F$ sont positives, et $i$
doit majorer la plus grande.} Cette propriété est stable par $F \mapsto pF$
(produit par le bi-isocristal de Tate) mais non par $F \mapsto p^{-1}F$ ; comme
on tient à inverser Tate, il n'est pas naturel d'imposer une condition de
croissance. Appelant \emph{effectifs} les $F$-isocristaux qui proviennent d'un
$\mathrm{IsBicr}(S,i)$, on a : tout $F$-isocristal est de la forme
$T^{-i} \otimes \mathcal{N}$ avec $\mathcal{N}$ effectif. « On n'a pas ajouté
plus de nouveaux objets qu'il n'était absolument nécessaire ! »

\pagerange{11}{15}
\subsection*{Exemple 2 : $S$ lisse sur un corps parfait (pages 11 à 15)}

\textbf{1.13.} Soit $S$ lisse sur un corps parfait $k$ de caractéristique $p$.
La catégorie $C$ ne change pas quand on remplace
$R = \operatorname{Spec}\mathbf{Z}$ par $\operatorname{Spec}W$, et un
cristal relatif à $W$ équivaut à un système compatible de cristaux relatifs aux
$W_n = W/p^{n+1}W$.\footnote{La lettre note ici que l'exemple 3 de 1.7 n'avait
pas été formulé avec la généralité voulue, celle d'une base formelle.} On peut
se localiser sur $S$ et le supposer affine ; il se relève alors en des $S_n$
lisses sur $W_n$, compatibles entre eux. Comme $S \to S_n$ est une immersion
fermée nilpotente, les cristaux sur $S$ relatifs à $W_n$ sont les cristaux sur
$S_n$ relatifs à $W_n$, et par 1.8 les objets de $S_n$ munis d'une
stratification relative à $W_n$. En résumé : \emph{un cristal sur $S$
s'identifie à un système compatible $(M_n)$ d'objets stratifiés sur les $S_n$
relativement aux $W_n$.}

Si $S$ se relève même en un schéma $X$ \emph{propre} et lisse sur $W$ —
hypothèse « évidemment bien restrictive » —, les théorèmes de comparaison
d'EGA~III, 4 et 5 donnent : \emph{les cristaux de modules cohérents sur $S$
forment une catégorie équivalente à celle des modules cohérents sur $X$ munis
d'une stratification relative à $W$}, catégorie qui se trouve donc
indépendante du relèvement. Sur la fibre générique $X_K$, schéma propre et
lisse sur le corps $K$ de caractéristique nulle, on obtient un foncteur de
restriction des cristaux cohérents sur $S$ vers les modules cohérents munis
d'une stratification relative à $K$,\footnote{La page dit « stratifiés
relativement à $X_K$ » ; il faut lire « à $K$ ».} c'est-à-dire, $X_K$ étant
lisse en caractéristique nulle, d'une connexion intégrable ; de tels modules
sont localement libres.\footnote{La lettre affirme plus généralement qu'un
module cohérent stratifié sur un schéma localement de type fini sur un corps
est localement libre ; on ne retient que le cas utilisé, $X_K$ lisse en
caractéristique nulle, qui est classique.}

Si de plus $X_K$ est géométriquement connexe et $K$ se plonge dans
$\mathbf{C}$, ces modules à connexion intégrable s'interprètent par les
représentations complexes du groupe fondamental de $X_{\mathbf{C}}^{\mathrm{an}}$.
Si ce groupe est trivial, tout module stratifié sur $X_K$ est trivial, et tout
cristal cohérent sur $S$ est \emph{isogène} à un cristal constant, image
inverse d'un cristal sur $k$ ; par rigidité, tout $F$-cristal cohérent, tout
bicristal de poids $i$, est isogène à un objet constant de même espèce. « On
voit donc là un principe d'approche transcendante pour l'étude des familles de
groupes formels en caractéristique $p$. » Le foncteur induit sur les
isocristaux cohérents est \emph{pleinement fidèle} ; est-ce une équivalence ?
La lettre trouve cela « un peu trop optimiste », et une note de sa main en
marge répond : « Effectivement, ça ne marche pas ! ».\footnote{Lecture de « ne
marche » incertaine. La condition cherchée sur un module stratifié, dit la
lettre, doit être « de nature locale en les points de $X$ qui sont maximaux
dans $S$ ». C'est ce qu'a confirmé la théorie ultérieure, pour les
isocristaux à puissances divisées : la condition est la convergence de la
connexion sur les disques résiduels des points fermés (isocristaux
convergents, A.~Ogus, 1984 ; cohomologie rigide de P.~Berthelot).}

Si $S$ n'est relevé que formellement, en un schéma formel $X$ sur
$\operatorname{Spf}W$, il reste vrai, de façon essentiellement triviale, que
\emph{les cristaux cohérents sur $S$ sont les modules cohérents sur $X$ munis
d'une stratification relative à $\operatorname{Spf}W$}. L'analogue de $X_K$ est
alors l'\emph{espace rigide-analytique} défini par $X$, avec ses modules
cohérents munis de stratifications au sens rigide-analytique (encore
localement libres), et l'on a encore un foncteur pleinement fidèle des
isocristaux cohérents sur $S$ vers ces modules. Il faut en déterminer l'image
essentielle, « car c'est vraiment ici la situation naturelle », et se demander
si ces modules ne se décrivent pas par un groupe « plus ou moins discret »
jouant le rôle du groupe fondamental. Conséquence immédiate de la
description : pour $S$ de type fini sur $k$, la catégorie des cristaux
cohérents est noethérienne, tout objet a un plus grand sous-objet de torsion,
et les objets sans torsion correspondent à des modules stratifiés sans
torsion, donc localement libres, sur $X$.

\textbf{1.14.} Il reste à décrire $\mathcal{M} \mapsto \mathcal{M}^{(p)}$. Un
carré commutatif $S' \to S$ au-dessus de $k' \to k$, relevé en $X' \to X$
au-dessus de $W' \to W$, explicite l'image inverse des cristaux par celle des
modules stratifiés. Pour le Frobenius, il s'agit de relever $S \to S^{(p/k)}$
en un morphisme de schémas formels
$f_X : X \to X$ compatible au Frobenius de $W$, ce qui est possible pour $S$
affine. Alors $\mathcal{M} \mapsto \mathcal{M}^{(p)}$ est l'image inverse
ordinaire par $f_X$ des modules stratifiés sur $X$ — resp., pour les
isocristaux, par le morphisme d'espaces rigides $X_K \to X_K$ qu'il induit,
semi-linéaire au-dessus du Frobenius de $K$. Le relèvement $f_X$ n'est pas
unique, mais le foncteur qu'il définit n'en dépend pas, à isomorphisme
canonique près.

\pagerange{15}{16}
\subsection*{Base quelconque (pages 15 et 16)}

\textbf{1.15.} Pour $S$ sans caractéristique fixée, soit $S_p = V(p)$ la fibre
de $S$ en $p$. Un bicristal de poids $i$ sur $S$ est un cristal de modules
$\mathcal{M}$ muni, pour chaque $p$, d'une structure $(F_p, V_p)$ de bicristal
de poids $i$ sur la restriction $\mathcal{M}_p$ à $S_p$. De même pour les
$F$-isocristaux, les isocristaux étant maintenant obtenus en tensorisant les
$\mathrm{Hom}$ par $\mathbf{Q}$. En caractéristique nulle la structure
supplémentaire est vide ; pour $S$ de type fini et dominant sur
$\operatorname{Spec}\mathbf{Z}$, elle est de nature arithmétique, les $F_p$
jouant le rôle des Frobenius, et en poids 1 le cristal « relie entre eux les
groupes formels en les diverses caractéristiques ».

\pagerange{16}{18}
\subsection*{Le site infinitésimal, et le test de De Rham (pages 16 à 18)}

\textbf{1.16. « Un retour en arrière. »} La définition 1.1 et la construction
de 1.3 sont « un petit peu canulées ». Il faut prendre pour $C$ la catégorie
des $R$-flèches $U \to T$, où $U$ parcourt les ouverts de $S$ et $U \to T$ est
une immersion fermée définie par un idéal nilpotent, et se borner aux
catégories fibrées $\mathcal{F}$ qui sont des \emph{champs} pour la topologie
de Zariski. On munit $C$ de la topologie de Zariski : une famille couvrante de
$(U \to T)$ est donnée par un recouvrement ouvert $(T_i)$ de $T$, avec
$U_i = U \cap T_i$. C'est le « \emph{site cristallogène} de $S$ sur $R$ », que
Grothendieck appellera ensuite \emph{site infinitésimal}.

Un faisceau $\mathcal{G}$ sur ce site est un système de faisceaux
$\mathcal{G}_{(U,T)}$ sur les $T$, avec, pour toute flèche
$(U \to T) \to (U' \to T')$, un morphisme de l'image inverse de
$\mathcal{G}_{(U',T')}$ sur $T$ vers $\mathcal{G}_{(U,T)}$, qui n'est pas
nécessairement un isomorphisme, mais l'est lorsque $T \to T'$ est une immersion
ouverte.\footnote{La parenthèse finale est un ajout de sa main au pied de la
page 16 ; « ouverte » y est de lecture incertaine, mais c'est ce que demande
la topologie de Zariski.} Les $\mathcal{O}_T$ forment un faisceau d'anneaux
$\mathcal{O}_C$, et \emph{les cristaux de modules quasi-cohérents sont
exactement les $\mathcal{O}_C$-modules quasi-cohérents} (localement conoyaux
de modules libres) ; les cristaux de présentation finie, les
$\mathcal{O}_C$-modules de présentation finie. Pour $S' \to S$, il n'y a pas de
morphisme naturel de sites, mais il y a un morphisme des \emph{topos}
cristallogènes $\mathrm{Topcr}_{S'/R} \to \mathrm{Topcr}_{S/R}$, construit
comme en 1.3, « qui, j'avoue, devrait être défini avec le plus grand soin ».

\textbf{1.17. Le test-clef.} On peut alors faire marcher la machine
cohomologique. Le test : \emph{si $R$ est le spectre d'un corps et $S$ lisse
sur $R$ (et propre en caractéristique $p > 0$), la cohomologie du site à
coefficients dans $\mathcal{O}_C$ doit être la cohomologie de De Rham de $S$
sur $R$} ; et pour $f : S' \to S$ propre et lisse, la valeur en $S$ de
$R^if_{\mathrm{cris}*}(\mathcal{O}_{C'})$ doit être $R^if_*(\Omega^\bullet_{S'/S})$,
la connexion de Gauss--Manin étant celle de la stratification canonique. Les
indices : l'existence de la connexion de Gauss--Manin, et « le tapis de
Washnitzer--Monsky ». À défaut, on dériverait dans la seule catégorie des
cristaux.

Une note au crayon dans la marge tranche : « marche en car.~0, et \emph{pas}
en car $p > 0$ ». C'est exact. En caractéristique nulle, la cohomologie du
site infinitésimal calcule la cohomologie de De Rham — c'est le théorème que
Grothendieck démontrera à l'IHÉS en décembre 1966. En caractéristique $p$,
elle ne la calcule pas : sur le site infinitésimal, $\mathcal{O}_C$ ne voit que
les fonctions annulées par \emph{tous} les opérateurs différentiels, y compris
les $\partial^n/n!$, alors que la cohomologie de De Rham ne voit que la
différentielle $d$. Le remède est celui qu'indique l'autre note au crayon : se
borner aux épaississements munis de puissances divisées. Le site qu'on obtient
est le site cristallin de Berthelot, et pour lui le test-clef est un
théorème.\footnote{Les deux notes au crayon de la page 17 sont postérieures à
la frappe ; la seconde, « (faut d'introduire des puiss.\ div.) », est lue en
grande partie par conjecture. La comparaison entre cohomologie cristalline et
cohomologie de De Rham, pour $S$ lisse et propre sur un corps parfait, est
dans Berthelot, \emph{op.\ cit.}, chapitre~V.}

\pagerange{18}{19}
\section*{La lettre, chapitre 2 : « La cohomologie de De Rham est un cristal » (pages 18 à 29)}

\subsection*{Trois indications (pages 18 et 19)}

\textbf{2.1.} Le titre est « un vœu pieux », mais la lettre le croit
« essentiellement correct ». Trois indications le soutiennent.

a) \emph{Le travail de Washnitzer et Monsky.} Il a deux défauts : il ne donne
que des invariants locaux sur une variété lisse en caractéristique $p$, ce qui
limite à la dimension 1 pour un invariant global ; et ses invariants sont des
espaces vectoriels sur $\operatorname{Frac}W$, c'est-à-dire « modulo
isogénie », les démonstrations d'invariance faisant intervenir des
dénominateurs. Il n'est donc pas exclu qu'il faille remplacer, dans le titre,
« cristal » par « isocristal » « (et "est" par "définit") ».\footnote{La
parenthèse est ajoutée à l'encre. La cohomologie de Monsky--Washnitzer est
publiée en 1968 (\emph{Formal cohomology I}, Annals of Math.\ 88) ; la lettre en
connaît les annonces.} Le test décisif reste 1.17 : marche-t-il seulement en
caractéristique nulle, ou en toute caractéristique ?

b) \emph{La connexion de Gauss--Manin.} Pour tout morphisme lisse
$f : X \to S$, il y a une connexion canonique (absolue) sur les
$R^if_*\Omega^\bullet_{X/S}$, et même sur le complexe
$Rf_*\Omega^\bullet_{X/S}$ dans la catégorie dérivée $D(\mathcal{O}_S)$. Sa
courbure est nulle, vérifiée par voie transcendante pour $S$ lisse sur une base
réduite de caractéristique nulle.\footnote{N.~Katz et T.~Oda en ont donné
en 1968 une construction algébrique, avec l'intégrabilité, par la
filtration de Leray du complexe de De Rham (\emph{J.~Math.\ Kyoto Univ.}~8).} En
caractéristique $p$, il n'y a pas d'espoir que cette connexion provienne
toujours d'une stratification : la cohomologie de De Rham d'une variété
affine y est « beaucoup trop grosse ». Il faut donc au moins supposer $f$
propre — condition nécessaire, dont 2.7 montrera qu'elle n'est pas
suffisante.

c) \emph{Le « papier bleu ».} Pour $S$ lisse sur $\mathbf{C}$ et $f : X \to S$
lisse et propre, la suite spectrale de Leray
\[
  E_2^{pq} = H^p\bigl(S^{\mathrm{an}}, R^qf^{\mathrm{an}}_*\mathbf{C}_{X^{\mathrm{an}}}\bigr)
  \Longrightarrow H^{p+q}(X^{\mathrm{an}}, \mathbf{C})
\]
a un terme initial et un aboutissement qui se décrivent algébriquement — par la
cohomologie de De Rham de $X$, et celle de $S$ à coefficients dans les
$R^qf_*\Omega^\bullet_{X/S}$ munis de leur connexion intégrable. On voudrait
une interprétation algébrique de la suite spectrale tout entière.\footnote{La
page écrit $\mathbf{C}_{S^{\mathrm{an}}}$ sous $f^{\mathrm{an}}_*$ ; il faut
lire $\mathbf{C}_{X^{\mathrm{an}}}$. Le « papier bleu sur De Rham » est
vraisemblablement l'article \emph{On the de Rham cohomology of algebraic
varieties} (Publ.\ Math.\ IHÉS 29, 1966) ; l'identification est nôtre. Une
interprétation algébrique de cette suite spectrale est précisément ce que
fournit la construction de Katz et Oda, par la filtration de Leray.}

\pagerange{19}{20}
\subsection*{Le formalisme souhaité (pages 19 et 20)}

\textbf{2.2.} « Dans l'éventualité optimiste où on n'aurait pas besoin de
s'isogéniser », les cristaux de modules forment une catégorie abélienne dont
on prend la catégorie dérivée : ses objets sont les \emph{complexes de
cristaux}, qui induisent sur chaque épaississement $T$ un complexe de
$\mathcal{O}_T$-modules ; un complexe de cristaux est parfait,
pseudo-cohérent\ldots\ s'il l'est sur chaque $T$. La théorie voulue associe à
tout morphisme lisse et propre $f : X \to S$ un complexe parfait de cristaux
absolus $\mathrm{DR}(f)$ sur $S$, la \emph{cohomologie de De Rham
cristalline} de $f$, avec :
\begin{itemize}
\item compatibilité à tout changement de base (au sens dérivé), et
fonctorialité contravariante en $X$ ;
\item une formule de Künneth
$\mathrm{DR}(f \times_S g) \simeq \mathrm{DR}(f) \overset{\mathbb{L}}{\otimes} \mathrm{DR}(g)$ ;
\item pour $f$ de dimension relative $d$, une dualité parfaite
\[
  \mathrm{DR}(f) \overset{\mathbb{L}}{\otimes} \mathrm{DR}(f) \longrightarrow T^d[-2d] ,
\]
où $T^d$ est le cristal de Tate de poids $2d$ ;\footnote{La page écrit
$\mathrm{DR}(f) \times \mathrm{DR}(f) \to T^d[2d]$, « où $[2d]$ indique qu'on
translate les degrés de $2d$ (attention au facteur 2 !) ». En notation
cohomologique, qui est celle d'aujourd'hui, l'accouplement $H^i \times H^{2d-i}
\to H^{2d}$ s'écrit avec $[-2d]$.}
\item un isomorphisme $\mathrm{DR}(f)(S) \simeq Rf_*(\Omega^\bullet_{X/S})$,
fonctoriel, compatible aux changements de base, à Künneth et à la dualité.
\end{itemize}
« Par prudence », la lettre ne dit rien du cas non propre, dont il faudrait
pourtant parler pour faire le lien avec Washnitzer--Monsky.\footnote{Rien de
ce formalisme n'est établi dans le dossier. Sa forme exacte, pour le site
cristallin à puissances divisées, est due à Berthelot (1974), avec Künneth et
la dualité de Poincaré.}

\pagerange{20}{21}
\subsection*{Le Frobenius sur la cohomologie (pages 20 et 21)}

\textbf{2.3.} Admettons la théorie de 2.2. Pour chaque premier $p$, notons
$\mathrm{DR}(f)_p$ la restriction à $S_p = V(p)$, qui par changement de base
est $\mathrm{DR}(f_p)$, avec $f_p : X_p \to S_p$. Par changement de base encore,
$\mathrm{DR}(f_p)^{(p)} = \mathrm{DR}(f_p^{(p)})$, où
$f_p^{(p)} : X_p^{(p/S)} \to S_p$ est le frobeniusé relatif. Le Frobenius
relatif $F_{X_p/S_p} : X_p \to X_p^{(p/S)}$ est un $S_p$-morphisme,\footnote{La
page écrit le morphisme de Frobenius « $X_p^{(p/S)} \to X$ ». La
contravariance de $\mathrm{DR}$ exige le sens $X_p \to X_p^{(p/S)}$, qui est
celui du Frobenius relatif.} et la fonctorialité contravariante de
$\mathrm{DR}$ donne
\[
  F_p : \mathrm{DR}(f)_p^{(p)} \longrightarrow \mathrm{DR}(f)_p .
\]
Il faut voir que c'est une isogénie, de sorte que
l'isocristal de $\mathrm{DR}(f)$ soit un $F$-isocristal. C'est ce que donne la
dualité de 2.2 : le cristal unité $T^d$ y est muni de sa structure de
bicristal, $F$ est compatible à l'accouplement, et en transposant $F$ on
obtient $V$ avec\footnote{La page écrit $FV = VF = p^{2d}$. Avec un accouplement à valeurs dans
$T^d$, sur lequel $F = p^d$, on a $\langle Fx, Fy\rangle = p^d \sigma\langle
x,y\rangle$, d'où $FV = VF = p^d$ pour $V$ transposé de $F$ ; pour une courbe
($d = 1$), on retrouve sur $H^1$ la relation $FV = p$ des modules de Dieudonné.
La conclusion de la lettre — $F$ est une isogénie — n'en dépend pas.}
\[
  FV = VF = p^d .
\]

Comme en caractéristique nulle, la formation des $\mathrm{DR}^i(f)$, une fois
passé aux isocristaux, devrait commuter aux changements de base, de sorte que
chaque $\mathrm{DR}^i(f)$ soit lui-même un $F$-isocristal. En rédigeant 1.10,
la lettre avait cru que $\mathrm{DR}^i(f)$ serait, sans passer aux isocristaux,
un bicristal de poids $i$ : « je m'étais canulé ». Pour définir $V$ sur
$\mathrm{DR}^i$ avec $FV = p^i$ en transposant le $F$ de $\mathrm{DR}^{2d-i}$,
il faudrait un \emph{isomorphisme} — pas seulement une isogénie —
$\mathrm{DR}^i(f) \simeq \mathrm{DR}^{2d-i}(f) \otimes T^{-(d-i)}$, comme en
donnerait une polarisation dans des cas « très exceptionnels ».

\pagerange{21}{24}
\subsection*{Schémas abéliens : les trois affirmations (pages 21 à 24)}

\textbf{2.4.} Pour un schéma abélien $A$ sur $S$, il semble raisonnable de
parler de bicristaux, et pas seulement de $F$-isocristaux : $\mathrm{DR}^1(A/S)$
porte les $F_p$ de 2.3, et des $V_p$ définis par fonctorialité à partir du
Verschiebung $A_p^{(p/S_p)} \to A_p$ ;\footnote{Construit en général dans les
exposés de Gabriel de SGA~3, que la lettre appelle « SGAD ».} comme
$\mathrm{DR}^1$ est multiplicatif en $A$ (Künneth) et que $FV = VF = p$ sur les
schémas abéliens en caractéristique $p$, $\mathrm{DR}^1(A)$ serait un
\emph{cristal de Dieudonné}, et sur un corps parfait un module de Dieudonné
sur $W(k)$ — qui doit être celui du groupe $p$-divisible de $A$.

Débarrassée de l'« hyperstructure axiomatico-conjecturale » de 2.2, la
construction se formule en trois affirmations.

\textbf{Première affirmation.} Soit $A$ un schéma abélien sur $S$ affine, et
$S \to T$ un épaississement nilpotent ; $A$ se prolonge en un schéma abélien
$B$ sur $T$, et $H^1_{\mathrm{DR}}(B/T)$ est localement libre de rang $2g$.
\emph{Ce module ne dépend, à isomorphisme canonique près, que de $A$ sur $S$
et de $S \to T$, et non du prolongement $B$.} On obtient ainsi un cristal
localement libre de rang $2g$, $\mathrm{DR}^1(A/S)$, additif en $A$, compatible
au changement de base, et défini sans hypothèse affine par recollement.

\textbf{Deuxième affirmation.} Sur $k$ parfait de caractéristique $p$, ce
cristal est un $W(k)$-module libre de rang $2g$, et, muni de $F$ et $V$, c'est
\emph{le module de Dieudonné $M$ de $A$} : pour tout anneau local artinien
$\Lambda$ de corps résiduel $k$ et tout prolongement $B$ de $A$ sur $\Lambda$,
$H^1_{\mathrm{DR}}(B/\Lambda) \simeq M \otimes_W \Lambda$, canoniquement et
fonctoriellement — la fonctorialité garantissant la compatibilité à $F$ et
$V$.

\textbf{2.5.} La lettre n'a pas vérifié ces deux affirmations, mais n'a « pas
le moindre doute » qu'elles soient correctes. Elles expliqueraient les
déformations de $A$ par son module de Dieudonné : les prolongements de $A$ sur
$\Lambda$ correspondraient aux quotients libres de rang $g$ de
$M \otimes_W \Lambda$ relevant le quotient canonique de
$M \otimes_W k = H^1_{\mathrm{DR}}(A)$.

\textbf{Troisième affirmation.} Pour un schéma abélien $f : A \to S$, la
cohomologie de De Rham porte la filtration de Hodge
\[
  0 \longrightarrow f_*\Omega^1_{A/S} \longrightarrow H^1_{\mathrm{DR}}(A/S)
  \longrightarrow R^1f_*\mathcal{O}_A \longrightarrow 0 .
\]
Pour tout prolongement $B$ de $A$ sur $T$, on a ainsi sur
$\mathcal{M}(T) = \mathrm{DR}^1(A/S)(T) = H^1_{\mathrm{DR}}(B/T)$ une filtration
qui prolonge celle de $\mathcal{M}(S) = H^1_{\mathrm{DR}}(A/S)$. \emph{Le
foncteur $B \mapsto (A, \text{filtration de } \mathcal{M}(T))$ est une
équivalence entre schémas abéliens sur $T$ et schémas abéliens sur $S$ munis
d'un relèvement de leur filtration de Hodge à $\mathcal{M}(T)$} (à quotients
localement libres).\footnote{La page écrit « prolongeant celle de
$\mathrm{DR}^1(A/A)(S)$ » : lire $\mathrm{DR}^1(A/S)(S)$. Elle note la
filtration par le signe tapé « Ø ».}

La filtration de Hodge n'est donc \emph{pas} de nature cristalline : elle
varie avec le relèvement, et c'est sa variation qui reflète les variations de
structure du motif. La comparaison avec le théorème de Serre--Tate — étendre
$A$ à $T$ en caractéristique $p$, c'est étendre son groupe $p$-divisible —
suggère une généralisation aux groupes formels, qui sera le chapitre 3.

\emph{Ce que valent ces trois affirmations.} Telles qu'elles sont énoncées, pour
\emph{tous} les épaississements nilpotents, elles sont fausses en
caractéristique $p$ : c'est exactement ce que la lettre elle-même démontre en
2.7, puisqu'elles fourniraient sur la cohomologie de De Rham des familles de
courbes elliptiques une stratification compatible au Frobenius. Elles sont
vraies, et ce sont des théorèmes, lorsqu'on se borne aux épaississements dont
l'idéal est muni de \emph{puissances divisées localement nilpotentes} : c'est
la théorie de Grothendieck--Messing. La première et la troisième sont le
théorème de Grothendieck--Messing (W.~Messing, \emph{The crystals associated to
Barsotti--Tate groups}, 1972 ; B.~Mazur et W.~Messing, \emph{Universal extensions
and one dimensional crystalline cohomology}, 1974) ; la deuxième, pour $\Lambda$
muni de telles puissances divisées (par exemple $\Lambda = W_n(k)$, $p \geq 3$),
est la comparaison entre $H^1$ cristallin et module de Dieudonné
contravariant (Mazur--Messing ; Berthelot, Breen et Messing, 1982).\footnote{Cette
mise au point est de la présente lecture, non de la page. Elle s'appuie sur
2.7 et 2.8, où la lettre établit elle-même l'impossibilité de la version sans
puissances divisées, et sur la note au crayon de la page 17.}

\pagerange{24}{25}
\subsection*{L'extension vectorielle universelle (pages 24 et 25)}

\textbf{2.6.} Paraphrasant sur une base quelconque une construction de Serre
« (c'est de lui que je l'ai apprise, du moins) », on trouve, pour tout schéma
abélien $A$ sur $S$, une \emph{extension vectorielle universelle}
\[
  0 \longrightarrow \omega_{A^\vee} \longrightarrow E(A) \longrightarrow A
  \longrightarrow 0 ,
\]
où $A^\vee$ est le schéma abélien dual et $\omega_{A^\vee}$, vu comme groupe
vectoriel, est le fibré vectoriel $\mathbb{V}(R^1f_*\mathcal{O}_A)$ dont les
sections sont les homomorphismes
$R^1f_*\mathcal{O}_A \to \mathcal{O}_S$.\footnote{La page écrit
$V(R^1f_*(\mathcal{O}_A)) = V(\underline{t}_{A\ })$, « où $A$ est le schéma
abélien dual de $A$ », avec un blanc laissé pour le signe du dual, qui n'a pas
été ajouté. Elle insiste sur la convention d'EGA : $\mathbb{V}(\mathcal{E})$ est
contravariant en $\mathcal{E}$.} Elle est universelle parmi les extensions de
$A$ par des groupes vectoriels, fonctorielle en $A$ et compatible au changement
de base. Son algèbre de Lie est une extension
\[
  0 \longrightarrow (R^1f_*\mathcal{O}_A)^\vee = \omega_{A^\vee}
  \longrightarrow \operatorname{Lie}E(A) \longrightarrow
  (f_*\Omega^1_{A/S})^\vee = \operatorname{Lie}A \longrightarrow 0 ,
\]
\emph{duale} de la filtration de Hodge : $\operatorname{Lie}E(A)$ s'identifie
au dual de $H^1_{\mathrm{DR}}(A/S)$. On a donc une construction de la suite de
Hodge par des extensions de groupes.

Les affirmations se précisent alors. La première : pour un prolongement $B$ de
$A$ à $T$, le groupe $E(B)$ ne dépend, à isomorphisme canonique — et même
unique — près, que de $A$ et de $S \to T$ ; on aurait un cristal en schémas en
groupes lisses $\mathcal{E}(A/S)$ avec $\mathcal{E}(A/S)(T) \simeq E(B)$.\footnote{La
page écrit $G(B/S)$ ; il faut lire $E(B/T)$, dans la notation de la page
$G(B/T)$, comme à la ligne suivante.} La troisième : relever la filtration de
Hodge à $\mathcal{M}(T)$, c'est relever le sous-groupe vectoriel canonique de
$\mathcal{E}(A/S)(S)$ en un sous-groupe vectoriel de $\mathcal{E}(A/S)(T)$, et
l'on retrouve $B$ comme quotient. C'est exactement la forme sous laquelle
Mazur et Messing démontreront ces énoncés, sur les épaississements à
puissances divisées nilpotentes.

\emph{Remarque sur Manin.} La connexion canonique de $E(A)$ est apparue en
cherchant à simplifier la construction par Manin de l'application
$A(S) \to \Gamma(S, \mathcal{O}_S)$ associée à une équation de
Picard--Fuchs.\footnote{L'« application de Manin », dans sa démonstration de la
conjecture de Mordell sur les corps de fonctions (1963). La page écrit
« Picart-Fuchs ».} Localement, toute section de $A$ se relève en une section
de $E(A)$ ; sa dérivée par la connexion est une section de
$\operatorname{Lie}E(A) \otimes \Omega^1_S$, définie modulo l'image de
$\omega_{A^\vee}$ — l'ambiguïté du relèvement —, et c'est un cocycle pour la
connexion de $\operatorname{Lie}E(A)$. Les équations de Picard--Fuchs servent
à en tirer une section de $\mathcal{O}_S$ ; du moins en caractéristique nulle,
où se place Manin.

\pagerange{26}{26}
\subsection*{Ce qui est vérifié (page 26)}

Les résultats positifs que la lettre dit avoir vérifiés :

a) \emph{En caractéristique nulle}, les affirmations 1 et 3, sous la forme
précisée par $E(A)$, sont vraies.

b) \emph{Sur une base quelconque}, $E(A)$ n'aurait pas d'automorphismes
infinitésimaux ; deux $E(B)$, $E(B')$ isomorphes le seraient donc de façon
unique, et l'hypothèse de 1 suffirait à produire un cristal en groupes, un
cristal de modules, et des stratifications canoniques. Cet énoncé est vrai en
caractéristique nulle et faux en caractéristique $p$, comme le dit une note au
crayon dans la marge, « faux ».\footnote{Lecture de « faux » douteuse.
Justification, qui n'est pas sur la page : les automorphismes infinitésimaux de
$E(A)$ correspondent aux homomorphismes $E(A) \to \mathbf{G}_a$, et la
suite exacte $0 \to \mathrm{Hom}(E(A), \mathbf{G}_a) \to \mathrm{Hom}(\omega_{A^\vee},
\mathbf{G}_a) \to \mathrm{Ext}^1(A, \mathbf{G}_a)$ montre que ce groupe est nul
lorsque les homomorphismes $\omega_{A^\vee} \to \mathbf{G}_a$ sont linéaires,
c'est-à-dire en caractéristique nulle ; en caractéristique $p$, les
homomorphismes $p$-polynomiaux forment un espace de dimension infinie, qui ne
peut s'injecter dans $\mathrm{Ext}^1(A, \mathbf{G}_a)$, de dimension finie.}

c) Les affirmations 1 à 3 sont vraies pour les relèvements \emph{d'ordre 1},
du moins lorsque $T$ est un schéma de nombres duaux
$D(\mathcal{J})$.\footnote{Une note au crayon, en biais, dit « à éclaircir,
\ldots\ peut être faux » (mot central illisible, dernier mot douteux). L'énoncé
au premier ordre est vrai : un idéal de carré nul porte les puissances
divisées triviales, nilpotentes, et l'on est dans le cadre de
Grothendieck--Messing.}

\pagerange{26}{27}
\subsection*{Le contre-exemple (pages 26 et 27)}

\textbf{2.7.} « Arrivé à ce point de mes brillantes conjectures, je m'aperçois
avec consternation qu'elles sont fausses telles quelles. » Soit $k$ un corps de
caractéristique $p > 0$. \emph{Il n'est pas possible de munir, pour tout
$k$-schéma $S$ de type fini de dimension $\leq 1$ et tout schéma elliptique $A$
sur $S$, le module $H^1_{\mathrm{DR}}(A/S)$ d'une stratification relative à
$k$, fonctorielle en $A$ et compatible aux changements de base.}\footnote{La
page ajoute « avec $S_{\text{rég}}$ régulier ». Une note à l'encre dans la
marge, entourée, porte « $\leq 0$ ! », le « 1 » étant barré d'un trait : la
correction resserre l'énoncé aux schémas de dimension nulle, et l'argument
suivant le permet, puisqu'il n'utilise que les voisinages infinitésimaux d'un
point. Le premier mot de la note n'est pas lu.}

L'ennui vient des courbes elliptiques de Hasse nul — les \emph{supersingulières}.
Les deux fonctorialités postulées rendent le Frobenius
$M^{(p)} \to M$, $M = H^1_{\mathrm{DR}}(A/S)$, compatible à la stratification.
En un point $s$ où $A_s$ est supersingulière, le Frobenius de
$H^1_{\mathrm{DR}}(A_s)$ est de carré nul ; un morphisme compatible à des
stratifications ayant un rang localement constant, on en déduit, en grimpant
sur les voisinages infinitésimaux de $s$, que la même propriété vaut aux points
voisins — ce qui est faux, par exemple pour la famille modulaire, où les points
supersinguliers sont isolés.

\emph{Conséquence.} Dans le titre du chapitre et dans 2.2, il faut prendre
partout des isocristaux pour garder une chance ; peut-être même la notion
d'isocristal est-elle encore trop restrictive, et faudra-t-il, dans la
description de 1.13, des stratifications en caractéristique nulle qui ne se
prolongent pas à tout le schéma formel sur $W$ — question qu'« une analyse
soigneuse du calcul-clef de Washnitzer--Monsky » devrait trancher.

L'argument est juste, et il porte exactement sur ce qu'il vise : les cristaux
\emph{sans} puissances divisées. Avec des puissances divisées, un cristal
n'impose plus une stratification mais seulement une connexion (à
$p$-courbure nilpotente), et un morphisme horizontal pour une connexion peut
changer de rang en caractéristique $p$ ; l'argument tombe, et le $H^1$
cristallin d'un schéma elliptique est bien un cristal, muni de son
Frobenius.\footnote{Ce paragraphe est de la présente lecture. Le $H^1$
cristallin d'un schéma abélien est un cristal à puissances divisées : c'est
le cas de dimension 1 de la théorie de Mazur et Messing (1974), et le
cristal de Dieudonné de son groupe $p$-divisible.}

\pagerange{27}{29}
\subsection*{Variétés abéliennes ordinaires (pages 27 à 29)}

\textbf{2.8.} Pour quels schémas abéliens les énoncés 1 et 3 valent-ils,
hors de la caractéristique nulle ? La lettre conjecture, pour $A$ variété
abélienne sur un corps $k$ de caractéristique $p$, que l'indépendance de $E(B)$
par rapport au relèvement $B$ sur tout anneau artinien de corps résiduel $k$
vaut \emph{si et seulement si $A$ est ordinaire}, c'est-à-dire a le nombre
maximum de points d'ordre $p$ sur une clôture algébrique de $k$. Le corps $k$
n'a pas à être parfait, et le groupe $p$-divisible de $A$ ne se décompose donc
pas nécessairement en étale et infinitésimal. On aurait alors, pour ces
variétés, la description des relèvements par relèvement d'une filtration d'un
module de Dieudonné, en accord avec la théorie de Serre des relèvements
canoniques (et des modules) des variétés abéliennes ordinaires.\footnote{Une note
au crayon, en biais et reliée à « si $A$ est une VA », dit : « faux, c'est
jamais vrai en car $p>0$ ! » (« c'est jamais » douteux), et un long trait
oblique au crayon traverse la fin de la page et la page suivante, comme pour
annuler. Ni la lettre ni le crayon ne démontrent rien ici ; nous ne tranchons
pas.}

Sur une base quelconque à fibres ordinaires, l'énoncé analogue est faux, même
pour les schémas elliptiques. Sur le schéma modulaire $S$, lisse sur
$\mathbf{Z}_p$, des courbes elliptiques à structure de niveau, une
stratification de $H^1_{\mathrm{DR}}$ relative à $\mathbf{Z}_p$ hors des
points supersinguliers — en nombre fini dans la fibre spéciale — s'étendrait à
tout $S$ « pour des raisons de pureté évidentes », serait compatible au
Frobenius, et contredirait 2.7. Il n'est pas exclu qu'en restant en
caractéristique $p$ l'énoncé vaille sur le lieu ordinaire, mais la
stratification obtenue ne s'étendrait pas à la courbe modulaire entière, et
il y aurait sur le lieu ordinaire une stratification naturelle en
caractéristique $p$ et une autre en caractéristique nulle « sans qu'elles
veuillent se recoller », ce qui « semble assez canularesque ».

Ces considérations font ressortir « la nature "infinie" de la notion de
stratification, par contraste avec celle de connexion, malgré les trompeuses
apparences de la caractéristique nulle ». Sur le $H^1_{\mathrm{DR}}$ de la
famille modulaire sur $\mathbf{Z}$, la stratification naturelle de la fibre
générique ne s'étend à aucun $S_U$, $U$ ouvert non vide de
$\operatorname{Spec}\mathbf{Z}$ — à la compatibilité au Frobenius près, qu'il
faudrait vérifier. Et c'est « assez moral » : le module stratifié joue le rôle
d'un faisceau $\ell$-adique, dont il partage la rigidité.

Même sur le lieu ordinaire, une stratification canonique du $H^1_{\mathrm{DR}}$
ne proviendrait pas d'un \emph{cristal} : elle se recollerait avec celle de la
caractéristique nulle et s'étendrait au schéma modulaire entier, points
supersinguliers compris, ce qui est absurde. Conclusion de la lettre : les
isocristaux, et les modules stratifiés, peuvent servir à l'étude des familles
de variétés abéliennes ; la notion de cristal elle-même « semble
irrémédiablement trop fine », sauf sur une base réduite à un point, ou au prix
de restrictions sérieuses qui figent la variation infinitésimale de $E(A)$ —
point de vue proche de celui de Serre, qui impose l'espace tangent et l'action
de la multiplication complexe. Cette conclusion vaut pour les cristaux de la
lettre ; pour les cristaux à puissances divisées, c'est la théorie de
Grothendieck--Messing qui la dément.

\pagerange{30}{31}
\section*{La lettre, chapitre 3 : remarques sur les groupes $p$-divisibles (pages 30 et 31)}

\textbf{3.1.} La construction de 2.6 se paraphrase pour un groupe
$p$-divisible $\Phi$ sur un schéma $S$ où $p$ est localement nilpotent. Les
homomorphismes de $\Phi$ dans le groupe formel additif sont nuls (en
particulier $\Phi$ n'a pas d'automorphismes infinitésimaux)\footnote{$\Phi$
est $p$-divisible et le groupe additif est tué par une puissance de $p$ sur
une telle base.} ; et le faisceau des $\mathrm{Ext}^1$ de $\Phi$ par le groupe
additif est localement libre de rang $g^* = \dim\Phi^*$, canoniquement
isomorphe à $\operatorname{Lie}(\Phi^*)$, où $\Phi^*$ est le groupe dual.\footnote{La
lettre l'énonce « sauf erreur », comme suggéré par les résultats de Lubin et
Tate sur les modules de groupes formels, et avoue ne pas avoir tiré au clair
ce qu'il faut entendre par « le groupe formel associé à $\mathbf{G}_a$ ».
L'énoncé est vrai, avec $\mathbf{G}_a$ lui-même (Messing, 1972).} D'où une
extension vectorielle universelle $E(\Phi)$ de $\Phi$ par le groupe vectoriel
$\omega_{\Phi^*}$, d'algèbre de Lie
\[
  (*) \qquad 0 \longrightarrow \omega_{\Phi^*} \longrightarrow
  \operatorname{Lie}E(\Phi) \longrightarrow \operatorname{Lie}\Phi
  \longrightarrow 0 ,
\]
fonctorielle en $\Phi$ et compatible au changement de base.

La lettre estime que $E(\Phi)$ varie « moins que $\Phi$ » mais n'est pas
indépendant des variations infinitésimales de $\Phi$ (sauf, probablement, au
premier ordre), et ne provient d'un cristal que si $\Phi$ est extension d'un
groupe ind-étale par un groupe torique, et $S$ réduit à un point. Cette
réserve vaut pour les cristaux de la lettre. Sur les épaississements à
puissances divisées nilpotentes, $E(\Phi)$ \emph{est} cristallin, pour tout
$\Phi$ : c'est le théorème de Messing, et $\operatorname{Lie}E(\Phi)$ est la
valeur du cristal de Dieudonné de $\Phi$. Dans le cas général, la lettre
propose d'étudier les variations de $\Phi$ à $E(\Phi)$ fixé par un cristal en
groupes.

\textbf{3.2.} L'extension $(*)$ subsiste par passage à la limite sur le
spectre d'un anneau noethérien $J$-adique séparé complet, par exemple un anneau
de valuation discrète complet $A$ d'inégales caractéristiques, de corps
résiduel $k$ parfait de caractéristique $p$. Soit $M$ le module de Dieudonné
(sur $W = W(k)$, de rang $h = g + g^*$) du groupe $p$-divisible réduit
$\Phi_k$.\footnote{La page écrit le rang « $g+g'$ » puis le type « $(g,g^*)$ » :
c'est le même nombre. Elle subordonne l'énoncé au cas où « le groupe formel
réduit a la structure triviale mentionnée plus haut dans 3.1 » ; cette
condition est obscure et l'énoncé qu'on donne n'en a pas besoin.} Alors la
donnée d'un relèvement de $\Phi_k$ en un groupe $p$-divisible sur $A$ équivaut
à celle d'un relèvement, en une filtration de $M \otimes_W A$ de type
$(g, g^*)$, de la filtration de $M \otimes_W k$ donnée par $(*)$ — \emph{lorsque
l'idéal maximal de $A$ porte des puissances divisées topologiquement
nilpotentes}, c'est-à-dire lorsque l'indice de ramification $e$ de $A$ est
$< p - 1$ (par exemple $A = W(k)$, $p \geq 3$).\footnote{Hypothèse ajoutée.
Sans elle, la valeur du cristal de Dieudonné sur $A$ n'est pas définie et
l'énoncé tombe ; la classification des groupes $p$-divisibles sur un anneau
très ramifié demande d'autres objets (modules de Breuil, de Kisin). Sous
l'hypothèse, c'est le théorème de Grothendieck--Messing, et pour
$A = W(k)$ la classification de J.-M.~Fontaine (\emph{Groupes $p$-divisibles
sur les corps locaux}, 1977).}

La lettre finit sur une question : \emph{« Il serait bien intéressant
d'essayer de déterminer, en termes d'une telle donnée, quel est le module de
Tate correspondant à la fibre générique de $\Phi$ ! »} — et sur une autre :
quelle quantité d'information l'extension $(*)$, qui remplace ici le $H^1$ de
De Rham, porte-t-elle pour des groupes plus généraux ? C'est la question que
Grothendieck posera au congrès de Nice en 1970 sous le nom de problème de la
« fonction mystérieuse », et à laquelle répond la théorie de Hodge $p$-adique
de Fontaine (représentations cristallines).\footnote{La lettre s'arrête là, au
milieu de la page 30 de sa pagination, sans formule finale ni signature.}

\pagerange{32}{36}
\section*{Réduction des courbes elliptiques et module de Tate $p$-adique (pages 32 à 36)}

Sous une chemise titrée de sa main « Réduction des courbes elliptiques et
$T_p(E)$ (cf.\ aussi livre de Serre "Abelian $l$-adic representations") »,
quatre pages de notes rapides, dont une grande part de la prose de liaison est
illisible ; on ne donne que ce que les formules et les mots lus
permettent d'établir.\footnote{Le livre de Serre est de 1968 : ces notes sont
au plus tôt de cette année-là. On les lit dans l'ordre 33, 34, 36, 35 : la page
36 s'arrête sur « Je conjecture », et la page 35 commence par « que cet élément
a une image non nulle dans $\mathbf{Z}_p$, donc que l'extension ne provient
pas d'une extension de groupes $p$-divisibles sur $A$ », qui ne peut suivre que
la discussion de la réduction multiplicative de la page 36 ; le cas c) de la
page 36, qui utilise $T_p(\mathbf{G}_m) = \mathbf{Z}_p(1)$ et la théorie de
Kummer, appartient d'ailleurs au cas 1º, de caractéristique nulle.}

Soit $A$ un anneau de valuation discrète complet, de corps des fractions $K$
et de corps résiduel $k$ de caractéristique $p > 0$, et $E_K$ une courbe
elliptique sur $K$. On veut décrire $T_p(E_{\bar K})$ comme module galoisien.

\textbf{1º) $K$ de caractéristique nulle} ; $T_p(E_{\bar K})$ est alors un
$\mathbf{Z}_p$-module libre de rang 2.

a) \emph{Bonne réduction ordinaire.} $E_K$ provient d'une courbe elliptique
$E$ sur $A$ et $E_k$ est ordinaire. Alors $T_p(E_K)$ est canoniquement une
extension
\[
  0 \longrightarrow L^\vee(1) \longrightarrow T_p(E_K) \longrightarrow L
  \longrightarrow 0 ,
\]
où $L$ est de rang 1 (le module de Tate de la partie étale) et
$L^\vee(1) = L^\vee \otimes T_p(\mathbf{G}_m)$ (celui de la partie
multiplicative, par l'accouplement de Weil). La classe de l'extension est dans
$H^1(K, L^{-2} \otimes T_p(\mathbf{G}_m))$, où $L^{-2} = (L^\vee)^{\otimes 2}$.
Si le corps résiduel est séparablement clos et $L$ trivial, de valeur $L_0$ en
$k$, la théorie de Kummer donne
\[
  H^1(K, \mathbf{Z}_p(1)) \otimes L_0^{-2} \simeq
  \Bigl(\varprojlim_n K^*/K^{*p^n}\Bigr) \otimes L_0^{-2}
  \simeq \mathrm{Hom}\bigl(L_0^{\otimes 2}, \widehat{K^*}\bigr) ,
\]
$\widehat{K^*}$ désignant le complété $p$-adique de $K^*$. La classe tombe
dans le sous-groupe $\mathrm{Hom}(L_0^{\otimes 2}, U^1)$, où $U^1 = 1 + \mathfrak{m}_A$
est le groupe des unités principales (« Einseinheiten ») : c'est que
l'extension provient d'une extension de groupes $p$-divisibles sur $A$.\footnote{C'est
le paramètre de Serre--Tate de la courbe ordinaire. La phrase qui suit, entre
crochets, et deux notes en marge, dont l'une semble dire que l'extension peut
être triviale, ne sont lues qu'en fragments.}

b) \emph{Bonne réduction supersingulière.} Il conjecture que la représentation
de $\pi_1(K)$ dans $T_p(E_{\bar K})$ est semi-simple, voire davantage sur la
taille de son image.\footnote{Il écrit « $\tfrac12$-simple » pour semi-simple ;
la fin de la phrase est illisible. La semi-simplicité est vraie : en réduction
supersingulière, l'inertie agit irréductiblement sur les points de
$p$-torsion (J.-P.~Serre, \emph{Propriétés galoisiennes des points d'ordre fini
des courbes elliptiques}, 1972), donc $T_p \otimes \mathbf{Q}_p$ est
irréductible.}

a$'$+b$'$) \emph{Bonne réduction potentielle} : l'invariant $j$ est entier, et
$E$ acquiert bonne réduction après extension finie de $K$. Les deux sous-cas,
ordinaire et supersingulier après extension, se ramènent aux précédents ; si la
conjecture de b) vaut, sa conclusion vaut aussi en b$'$). Il conjecture en
outre, pour $j$ entier, que \emph{$E_K$ a bonne réduction si et seulement si
son groupe $p$-divisible a bonne réduction}, c'est-à-dire s'étend en un groupe
$p$-divisible sur $A$.\footnote{L'énoncé est vrai, et sans hypothèse sur $j$,
pour les variétés abéliennes sur $K$ de caractéristique nulle : bonne réduction
équivaut au caractère cristallin de $T_p \otimes \mathbf{Q}_p$ (R.~Coleman et
A.~Iovita, 1999 ; C.~Breuil, 2000), qui équivaut à l'extension du groupe
$p$-divisible (Breuil, Kisin).}

c) \emph{L'invariant $j$ n'est pas entier.} Le modèle de Néron $E$ a alors
une composante neutre $E^0_k$ forme de $\mathbf{G}_m$ ; quitte à faire une
extension finie, $E[p^\nu]$ contient un sous-groupe canonique forme de
$\mu_{p^\nu}$, et $T_p(E_K)$ un sous-module forme tordue de
$T_p(\mathbf{G}_m)$. Si $k$ est séparablement clos, on a une extension
\[
  0 \longrightarrow \mathbf{Z}_p(1) \longrightarrow T_p(E_K) \longrightarrow
  \mathbf{Z}_p \longrightarrow 0
\]
(le quotient étant $\mathbf{Z}_p$ car $\det T_p(E_K) \simeq \mathbf{Z}_p(1)$),
déterminée par un élément canonique de $\widehat{K^*}$, lui-même extension de
$\mathbf{Z}_p$ (venant de $K^*/A^*$) par $U^1$. Il conjecture que cet élément
a une image \emph{non nulle} dans $\mathbf{Z}_p$, de sorte que l'extension ne
provient pas d'une extension de groupes $p$-divisibles sur $A$, et demande :
\emph{quel est cet élément de $\mathbf{Z}_p$ ? est-ce toujours un entier, un
élément $> 0$ ?}\footnote{Les deux réponses sont oui. Par l'uniformisation de
Tate, $E_K$ devient, après extension finie, $\mathbf{G}_m/q^{\mathbf{Z}}$ avec
$v(q) = -v(j) > 0$, et l'élément de $\mathbf{Z}_p$ est l'entier $v(q)$. Les
crochets de la page (« de façon canonique au signe près ») sont gardés dans le
sens : l'identification est canonique au signe près.}

\textbf{2º) $K$ de caractéristique $p$.} Il faut alors étudier $T_p(E_{K_s})$,
qui n'a plus que le rang 0 ou 1, et que les seuls modules galoisiens ne
décrivent plus. En bonne réduction ordinaire, $T_p(E_K)$ est encore une
extension de $L_K$ par $L_K^\vee \otimes T_p(\mathbf{G}_m)$, donnée par un
élément de $L^{-2} \otimes U^1$\ldots\ La note s'arrête là.\footnote{En
caractéristique $p$, $T_p(\mathbf{G}_m)(K_s) = 0$, et la formule doit
s'entendre pour les groupes $p$-divisibles plutôt que pour les points ; la
page s'interrompt sur un « $\otimes$ ».}

\pagerange{37}{38}
\section*{Une feuille de Tate (pages 37 et 38)}

Sous une chemise de sa main, « Gribouillis de Tate sur sa théorie », une feuille
d'une autre main, en anglais, à l'encre très pâlie : vraisemblablement des notes
de Tate. On y lit l'hypothèse « $L/K$ Galois », l'anneau des entiers $R_L$, le
groupe $\mathcal{G} = \mathrm{Gal}(L/K)$, un théorème reliant
$\widehat{L}^{\mathcal{G}} = \widehat{K}$ à la nullité de
$H^1(\mathcal{G}, R_L)$, et des suites exactes. Ce n'est pas de Grothendieck,
et le dossier n'y porte aucune intervention de sa main ; on ne la reprend
pas.\footnote{Le sujet — invariants galoisiens du complété d'une extension et
cohomologie galoisienne des entiers — est celui de la théorie de Tate des
pages suivantes (J.~Tate, \emph{$p$-divisible groups}, conférence de Driebergen,
1966, publiée en 1967).}

\pagerange{39}{44}
\section*{Questions de théorie de Tate--Hodge (pages 39 à 44)}

Sous la chemise « Questions diverses en théorie de Tate--Hodge », trois pages
écrites au dos d'un tapuscrit étranger au dossier. Le cadre : $K$ corps local
de caractéristique nulle, de corps résiduel de caractéristique $p$,
$\pi = \mathrm{Gal}(\bar K/K)$, et $\mathbf{C}$ le complété de $\bar K$, sur
lequel $\pi$ agit ; $\mathbf{C}(i)$ est le tordu de Tate.\footnote{Les pages
écrivent $\mathbf{C}[i]$ pour le tordu de Tate, et
$\mathbf{Q}_\ell$, $T_\ell$ là où le contexte impose $\ell = p$ ; on écrit
$\mathbf{C}(i)$ et $\mathbf{Q}_p$. Le crayon des archivistes numérote la
chemise et la première page 50 et 51.}

\textbf{Les questions (page 40).}
\begin{itemize}
\item[a)] Y a-t-il un isomorphisme entre cohomologie de De Rham et cohomologie
de Hodge après extension des scalaires à $\mathbf{C}$, fonctoriel, multiplicatif
et compatible aux filtrations, une fois choisi un générateur
$\varepsilon$ de $T_p(\mathbf{C}^*) = \mathbf{Z}_p(1)$ ?\footnote{Lecture
douteuse des premiers mots (« Théorème ») et de deux mots au-dessus de
« filtrations ». La réponse moderne passe par l'anneau $B_{\mathrm{dR}}$ de
Fontaine, où le choix de $\varepsilon$ fixe l'uniformisante $t$, et par le
théorème de comparaison de G.~Faltings (1988--89).}
\item[b)] Tout $\mathbf{Q}_p[\pi]$-module simple de dimension finie est-il
contenu dans un $\mathbf{C}(i)$ ?\footnote{Lecture douteuse de toute la phrase
et de la suivante. Telle quelle, la réponse est non : pour un caractère
$\chi = \chi_{\mathrm{cycl}}^a$ avec $a \in \mathbf{Z}_p \setminus \mathbf{Z}$,
le théorème de Tate ($H^0(\pi, \mathbf{C}(\psi)) = 0$ dès que $\psi$ est d'image
infinie sur l'inertie) montre que $\mathbf{Q}_p(\chi)$ ne se plonge dans aucun
$\mathbf{C}(i)$ ; la théorie de Sen (1980) dit quels modules se plongent.}
\item[c)] Les $M = H^i(X_{\mathbf{C}}, \mathbf{Z}_p) \otimes \mathbf{Q}_p$
satisfont-ils d'autres conditions que celles qui résultent de la conjecture de
Tate ?
\item[d)] \emph{Construire} les isomorphismes de Tate
$H^{p,q}(X)_K \simeq H^q(X, \Omega^p_X)$.
\item[e)] Quel lien avec les connexions canoniques ? f) Avec
Washnitzer--Monsky ?
\end{itemize}

\textbf{Trois structures sur un module de Hodge--Tate (page 42).} Pour un
groupe $p$-divisible sur $K$, on a sur $V = M \otimes_{\mathbf{Z}_p} \mathbf{C}$ :
\begin{enumerate}
\item une décomposition $V = V^{10} \oplus V^{01}$ en $\mathbf{C}$-espaces ;
\item une action semi-linéaire de $\pi$ ;
\item un réseau $M$ libre de type fini sur $\mathbf{Z}_p$.
\end{enumerate}
Relations : 1 se déduit de 2, par
\[
  V^{10} = V^\pi \otimes_K \mathbf{C}, \qquad
  V^{01} = \bigl(V \otimes \mathbf{C}(1)\bigr)^\pi \otimes_K \mathbf{C}(-1)
\]
(« $2 \Rightarrow 1$ ») ;\footnote{C'est le théorème de Tate (1967) : l'action
semi-linéaire de Galois détermine la décomposition de Hodge--Tate. Les indices
et les tordus sont de lecture en partie douteuse ; une ligne biffée proposait
une autre écriture.} 2 et 3 ensemble, c'est un réseau stable par $\pi$,
c'est-à-dire le groupe $p$-divisible sur $K$ ; entre 3 et 1 : « ??? ». Le corps
$K$ n'intervient que dans 2 ; 1 et 3 sont intrinsèques au groupe $p$-divisible
sur $\mathbf{C}$. Il propose enfin de regarder \emph{le plus petit groupe
algébrique $G$ d'automorphismes de $M \otimes \mathbf{Q}_p$ tel que
$G_{\mathbf{C}}$ contienne les opérations $\operatorname{diag}(\lambda, \xi)$
sur la décomposition $V^{10} \oplus V^{01}$}.\footnote{La matrice est de lecture
douteuse. C'est l'analogue $p$-adique du groupe de Mumford--Tate ; que l'algèbre
de Lie de l'image de l'inertie soit la plus petite $\mathbf{Q}_p$-algèbre de Lie
algébrique dont l'extension à $\mathbf{C}$ contienne l'opérateur de
Hodge--Tate est le théorème de S.~Sen (Annals of Math., 1973), conjecturé par
Serre.}

\textbf{La conjecture de Hodge--Tate (page 44).} Pour $X$ propre et lisse sur
$K$, notant $H^i_B(X_{\mathbf{C}}, \mathbf{C}) =
H^i_{\text{ét}}(X_{\bar K}, \mathbf{Q}_p) \otimes \mathbf{C}$ et, pour $p + q = i$,
$H^{p,q}_B(X)_K = \bigl(H^i_B(X_{\mathbf{C}}, \mathbf{C})(p)\bigr)^\pi$ :

\emph{Conjecture.} 1) On a une décomposition
\[
  \boxed{\;H^i_B(X_{\mathbf{C}}, \mathbf{C}) \simeq \bigoplus_{p+q=i}
  H^{p,q}_B(X)_K \otimes_K \mathbf{C}(-p)\;}
\]
2) On a un isomorphisme canonique
\[
  \boxed{\;\rho^{pq} : H^{p,q}_B(X)_K \simeq H^q(X, \Omega^p_X)\;}
\]
3) Pour la classe $P_B(2i) \in H^{2i}(X_{\mathbf{C}}, \mathbf{Z}_p(i))^\pi
\subset H^{i,i}_B(X)_K$ et la classe $P_{\mathrm{DR}}(2i) \in H^{i,i}(X)$, on
veut $\rho^{ii}(P_B(2i)) = P_{\mathrm{DR}}(2i)$.\footnote{La page ne dit pas
d'où viennent les classes $P(2i)$ : classes d'un cycle de codimension $i$, ou
puissance $i$-ième d'une classe de polarisation. La décomposition 1)--2) est le
théorème de Faltings (1988), conjecturé par Tate en 1967 ; la compatibilité aux
classes de cycles est établie avec la comparaison cristalline et de De Rham
(Faltings, Tsuji).}

\pagerange{46}{48}
\section*{La question d'isogénie de Barsotti (pages 46 à 50)}

La chemise : « Question d'isogénéité de Barsotti ; impossibilité de plonger un
groupe formel $p$-divisible dans une v.A.\ sur un $k$ donné (non alg.\ clos) ».

\textbf{1) La question de Barsotti.} Soient $A$, $B$ deux variétés abéliennes
sur $k$ algébriquement clos de caractéristique $p$, « hypersingulières »,
c'est-à-dire sans point d'ordre $p$. Si leurs groupes formels sont isogènes,
$A$ et $B$ sont-elles isogènes ? En dimension 1, c'est dans Deuring. Manin ignore
la réponse en dimension 2, où les groupes formels sont isogènes à
$G_{1,1} + G_{1,1}$ : la question est alors de savoir si une surface abélienne
hypersingulière est isogène au carré d'une courbe elliptique
supersingulière.\footnote{« Hypersingulière » est la lecture de la page ;
« sans points d'ordre $p$ » (rang $p$ nul) équivaut à supersingulière en
dimension $\leq 2$, non au-delà. Pour les variétés supersingulières, la
réponse est oui : toute variété abélienne supersingulière est isogène à une
puissance d'une courbe elliptique supersingulière (F.~Oort, \emph{Subvarieties
of moduli spaces}, 1974, à partir du théorème de Tate sur les corps finis).
La page note que Manin l'affirme en caractéristique 2.}

Sans l'hypothèse hypersingulière, la réponse est négative en dimension 2, et la
page le démontre. D'après Manin, les surfaces abéliennes polarisées de groupe
formel $\widehat{\mathbf{G}}_m + G_{1,1}$ forment une sous-famille $N$ de
dimension 2 de la famille modulaire $M_2$. Parmi elles, celles qui ne sont pas
simples sont isogènes à un produit $E \times F$, $E$ ordinaire et $F$
supersingulière ; sur un corps algébriquement clos elles sont même
isomorphes à de tels produits, et ces produits forment une famille de dimension
1 au plus, image de l'ouvert ordinaire de la droite modulaire par
$E \mapsto E \times F$, $F$ parcourant un nombre fini de classes.\footnote{La
page écrit « de dim 2 au plus » ; c'est 1 qu'il faut lire, ce que confirme
l'argument lui-même (une courbe modulaire, un nombre fini de $F$). Elle dit
aussi « les courbes isogènes à $E\times F$ » pour « les variétés ». Un
passage barré de croix suit, lu en fragments.} D'où un ouvert non vide
$U \subset N$ formé de variétés simples. Or les variétés isogènes à une
variété simple $A$ de ce type forment une famille \emph{discrète} : si $A$ est
définie sur $k_0$, toute $A'$ isogène l'est sur une extension finie de
$k_0$.\footnote{Justification, qui n'est pas sur la page : les sous-groupes
finis de $A[p^\infty] = \widehat{\mathbf{G}}_m \times G_{1,1} \times
\mathbf{Q}_p/\mathbf{Z}_p$ d'un ordre donné sont en nombre fini, car la partie
locale-locale $G_{1,1}$ n'a qu'un sous-groupe de chaque ordre, et il n'y a pas
d'homomorphismes non nuls entre les trois parties.} Deux points de $U$ de degrés
de transcendance différents sur $\mathbf{F}_p$ (0, 1 ou 2) donnent donc des
variétés simples \emph{non isogènes} de même groupe formel
$\widehat{\mathbf{G}}_m \times G_{1,1}$.

\pagerange{49}{50}
\subsection*{Un groupe formel qui ne se plonge dans aucune variété abélienne (pages 49 et 50)}

\textbf{2)} Sur un corps $k$, même fini, un groupe formel — même une forme de
$G_{1,1}$ — ne se réalise pas en général comme sous-groupe formel du groupe
formel d'une variété abélienne définie sur $k$, même après extension finie de
$k$. Soit $k$ fini à $q$ éléments et $E$ une courbe elliptique supersingulière
sur $k$ dont tous les endomorphismes sont définis sur $k$, de groupe formel
$G_0 = \widehat{E}$. Alors $\pi_1(k)$ agit trivialement sur $\mathrm{End}(E)$,
ordre maximal d'une algèbre de quaternions sur $\mathbf{Q}$, et sur
$\mathrm{End}(G_0) \simeq \mathrm{End}(E) \otimes \mathbf{Z}_p = \mathcal{A}$,
ordre maximal d'une algèbre de quaternions sur $\mathbf{Q}_p$. Tordons $G_0$
par une unité $\theta$ de $\mathcal{A}$ : le nouveau Frobenius est
$\varphi' = \theta\varphi$. Si le tordu se plongeait dans le groupe formel
d'une variété abélienne $A$ sur $k$, les valeurs propres de $\varphi'$ seraient,
par Weil, des entiers algébriques de valeurs absolues $q^{1/2}$, et
$\det\varphi' = \det\varphi \cdot \det\theta$ serait algébrique, donc
$\det\theta$ aussi. Or il est facile de trouver une unité $\theta$ de
$\mathcal{A}$ dont la norme réduite $\det\theta$ n'est pas algébrique.\footnote{La
norme réduite $\mathcal{A}^* \to \mathbf{Z}_p^*$ est surjective, et
$\mathbf{Z}_p^*$ contient des éléments transcendants ; une extension finie de
$k$ remplace $\varphi'$ par une puissance, ce qui ne change rien. Un bloc de
quatre lignes biffé, page 49, est lu en fragments. « $\det$ » est la norme
réduite.}

\pagerange{51}{53}
\section*{Déformations d'extensions de groupes de Barsotti--Tate, et schémas modulaires (pages 51 à 56)}

La chemise : « [Prolongements infinitésimaux d'extension de groupes de B.T.\ :]
Conséquences pour structure infinitésimale des schémas modulaires ». La première
page commence au milieu d'une phrase ; la feuille qui la précédait n'est pas
dans le dossier à cette place.\footnote{Les pages 52 à 56 sont d'une écriture
rapide et pâle ; une grande part de la prose de liaison est illisible. On
restitue l'argument par ses formules, et on le dit là où la restitution est
nôtre.}

\textbf{Le point de départ (page 52).} Pour une extension $E$ d'un groupe de
Barsotti--Tate $M$ par un autre $N$, sur les voisinages infinitésimaux
$S^{(n)}$ de $S$ dans une base $R$, les morphismes classifiants
$S^{(n)} \to G$ vers un groupe formel $G$, nuls sur $S^{(0)}$, se recollent en
un morphisme $S^{(\infty)} \to G$ dont la donnée équivaut à celle de
l'extension sur le voisinage infinitésimal de $S$ ; sa formation commute au
changement de base. D'où l'idée de la lire dans une situation
« \emph{modulaire} ».

\textbf{Modules formels d'extensions (pages 52 et 53).} Soit
$R = \operatorname{Spec}\Lambda$, $\Lambda$ local complet de caractéristique
résiduelle $p$ (par exemple un anneau de Cohen), de corps résiduel $k$, et
$M_0$, $N_0$ des groupes de Barsotti--Tate sur $k$, $M_0$ étale et $N_0$ de type
multiplicatif.\footnote{« $M_0$ étale » est une hypothèse que la page ne
formule pas mais qu'elle utilise, puisqu'elle écrit $T_p(M_0)$.} On considère
la variété formelle modulaire $\mathcal{M}_\xi$ sur $\operatorname{Spf}\Lambda$
des déformations d'une extension $\xi$ de $M_0$ par $N_0$. Alors :
\begin{enumerate}
\item pour l'extension triviale $\xi_0$, $\mathcal{M}_{\xi_0}$ est un groupe
formel lisse, et même un \emph{tore formel} — donc déterminé par sa réduction —,
de réduction $T_p(M_0)^\vee \otimes N_0$ ;
\item $\mathcal{M}_\xi$ est un \emph{torseur} sous $\mathcal{M}_{\xi_0}$.
\end{enumerate}
Si l'on suppose seulement $N_0$ connexe, une extension $E_0$ détermine $M_0$,
$N_0$ et la structure d'extension, et le morphisme naturel
$X = \mathcal{M}_{E_0} \to \mathcal{M}_{N_0} = Y$ fait de $X$ un torseur sous
$\mathcal{N}_Y \otimes T_p(E_0)^\vee_Y$ ; la structure de
$\mathcal{M}_{E_0}$ est connue dès que l'est celle de
$\mathcal{M}_{N_0}$.\footnote{Comme $N_0$ est connexe, les points de $E_0$ sur
une clôture algébrique de $k$ sont ceux de $M_0$, et $T_p(E_0) = T_p(M_0)$. Pour
$N_0$ de type multiplicatif, $Y$ est réduit à un point.}

Pour $M = \mathbf{Q}_p/\mathbf{Z}_p$ et $N = \mu_{p^\infty}$, le tore de 1 est
$\widehat{\mathbf{G}}_m$ : c'est la coordonnée $q$ de Serre--Tate.

\pagerange{54}{56}
\subsection*{Couples compatibles et variétés abéliennes polarisées (pages 54 à 56)}

Soient des homomorphismes $u : M \to M'$ et $v : N \to N'$, et des extensions
$\xi$ de $M$ par $N$ et $\xi'$ de $M'$ par $N'$, triviales sur $S_0$. Il existe
un morphisme d'extensions compatible à $u$ et $v$ si et seulement si
$v_*(\xi) = u^*(\xi')$ dans $\mathrm{Ext}^1(M, N')$. Or $v_*$ est
« représenté » par l'homomorphisme de tores
$\alpha : T_p(M)^\vee \otimes N \to T_p(M)^\vee \otimes N'$ défini par $v$, et
$u^*$ par $\beta : T_p(M')^\vee \otimes N' \to T_p(M)^\vee \otimes N'$ défini par
$u$. Les couples $(\xi, \xi')$ compatibles sont donc représentés,
infinitésimalement, par le \emph{noyau} de
\[
  \alpha - \beta : G \times G' \longrightarrow H .
\]
Pour $N$, $N'$ de type multiplicatif, la variété formelle des modules est
ainsi un groupe formel de type multiplicatif sur $\Lambda$, noyau d'un
homomorphisme de tores formels.

\emph{Application aux variétés abéliennes polarisées ordinaires.} Pour une
variété abélienne ordinaire, le groupe $p$-divisible est extension d'un groupe
étale par un groupe de type multiplicatif, et une polarisation fournit les
homomorphismes $u$, $v$. Si $\Lambda$ est un corps de caractéristique $p$,
l'anneau local du schéma modulaire des variétés abéliennes polarisées
ordinaires (polarisation donnée, rigidification de Jacobi d'échelon donné) en
un point est donc une \emph{intersection complète} — en particulier de
Cohen--Macaulay et de Gorenstein —, \emph{localement irréductible}, de
réduit lisse sur $k$. Comme la dimension de ces schémas modulaires est
$\tfrac{g(g+1)}{2}$ en chaque point, on conclut : \emph{si la polarisation est
de degré premier à $p$, le schéma modulaire des variétés abéliennes ordinaires
est lisse en caractéristique $p$.}\footnote{Une note en marge, en travers, dit à
peu près : « dimensionnalité ridicule, on le voit bien directement en
car.~$p$ ! » (lecture douteuse). Plusieurs mots de l'argument sont illisibles ;
la restitution est nôtre, et elle est conforme au résultat connu : pour $p$ ne
divisant pas le degré, l'espace des déformations polarisées d'une variété
ordinaire est le sous-tore des paramètres de Serre--Tate symétriques, de
dimension $g(g+1)/2$ (N.~Katz, \emph{Serre--Tate local moduli}, 1981). La
page insiste : il ne s'agit que de variétés ordinaires.}

\pagerange{57}{59}
\section*{Compléments sur les groupes de Witt (pages 58 et 59)}

Sous ce titre, souligné de sa main, des énoncés en vue d'une théorie de
Dieudonné des groupes unipotents par les vecteurs de Witt. Toute la page 59 est
barrée de deux longs traits obliques : ses énoncés sont des brouillons retirés,
qu'on rapporte comme tels.\footnote{Ces pages citent Demazure--Gabriel,
\emph{Groupes algébriques} (1970), « II~§3 4.6 » et « V~§1 2.2 » : elles sont au
plus tôt de 1970, ce qui s'accorde avec la datation du catalogue. La page 57
est blanche.}

\textbf{Lemme 1.} Soit $S = \operatorname{Spec}A$ de caractéristique $p$. Le
faisceau $\underline{\mathrm{Ext}}^1(\mathbf{G}_a, \mathbf{G}_a)$ est le
$A$-module (à gauche) libre de base les $\mathbf{F}^i e$, c'est-à-dire le
$A[\mathbf{F}]$-module libre de base $e$, où $e$ est la classe de $W_2$ comme
extension de $\mathbf{G}_a$ par $\mathbf{G}_a$. \emph{Démonstration} : on se
ramène à $A = \mathbf{F}_p$, et c'est dans Demazure--Gabriel.

\textbf{Lemme 2.} Soit $H$ un sous-groupe de $\mathbf{G}_a$ tel que
$\mathbf{G}_a/H$ soit localement isomorphe à $\mathbf{G}_a$.\footnote{L'énoncé
de l'hypothèse est en partie biffé et illisible ; on garde la condition lisible.
Exemples : $H = \alpha_{p^n}$, ou $H = \mathbf{Z}/p$, noyau de $x \mapsto x^p -
x$.} Alors :
\begin{itemize}
\item[a)] $\underline{\mathrm{Ext}}^1(H, W_n) \to \underline{\mathrm{Ext}}^1(H, W_1)$
est bijectif ; a$_1$) $\mathrm{Ext}^1(H, W_{n-1}) \to \mathrm{Ext}^1(H, W_n)$
est nul, i.e.\ $\mathrm{Hom}(H, \mathbf{G}_a) \to \mathrm{Ext}^1(H, W_{n-1})$ est
surjectif ; a$_2$) $\mathrm{Ext}^1(H, W_1) \to \mathrm{Ext}^1(H, W_{n-1})$ est
nul ;
\item[b)] $\mathrm{Ext}^1(\mathbf{G}_a, W_n) \to \mathrm{Ext}^1(H, W_n)$ est
surjectif ;
\item[c)] $\underline{\mathrm{Ext}}^1(\mathbf{G}_a, W_n)$ est un
$A[\mathbf{F}]$-module à droite engendré par la classe de $W_{n+1}$.
\end{itemize}
\emph{Démonstration} : pour $n = 1$, a) est trivial, b) résulte de
$\mathrm{Ext}^2(\mathbf{G}_a/H, W_1) = 0$, conséquence par dévissage de
$\mathrm{Ext}^2(\mathbf{G}_a, \mathbf{G}_a) = 0$, et c) est le lemme 1 ; en
général, par récurrence sur $n$, comme dans Demazure--Gabriel.\footnote{La
flèche de a) porte « $R^{n-1}$ » (restriction) ; une ligne biffée et un petit
schéma entre a) et a$_1$) sont de lecture douteuse.}

\textbf{Corollaire 3} (page barrée). Pour $H$ comme ci-dessus,
$\underline{\mathrm{Ext}}^1(H, \underline{W}) = 0$, où
$\underline{W} = \varinjlim W_n$. La question en marge — a-t-on
$\mathrm{Ext}^1(H, \underline{W}) = 0$ globalement, i.e.\
$H^1(S, \underline{\mathrm{Hom}}(H, \underline{W})) = 0$ ? — reçoit la réponse :
\emph{non}, déjà pour $H = \alpha_p$ si $S$ n'est pas parfait.

\textbf{Proposition 4} (page barrée). Soit $G$ un groupe algébrique fini (par
exemple sur un corps parfait) annulé par $V^n$. Alors
$\underline{\mathrm{Hom}}_{\mathrm{gr}}(G, W_n)$ est un schéma en
groupes,\footnote{La suite de a) — une description de ce schéma en groupes en
fonction du rang $p^r$ de $G$ — est en partie biffée et de lecture douteuse.} et
$\underline{\mathrm{Ext}}^1_{\mathrm{gr}}(G, \underline{W}) = 0$. Le b) vient du
corollaire 3 par dévissage en $\alpha_p$ ; le a) par récurrence sur la longueur
d'une suite de composition $0 \to G' \to G \to \alpha_p \to 0$, qui donne
\[
  0 \to \underline{\mathrm{Hom}}(\alpha_p, \underline{W}) \to
  \underline{\mathrm{Hom}}(G, \underline{W}) \to
  \underline{\mathrm{Hom}}(G', \underline{W}) \to 0 ,
\]
avec $\underline{\mathrm{Hom}}(\alpha_p, W_n) \simeq \mathbf{G}_a$. « On gagne ! »

\textbf{5} (page barrée). Posons $\mathbb{D}(G) = \underline{\mathrm{Hom}}(G,
\underline{W})$. C'est un module sur les vecteurs de Witt, muni de $\mathbf{F}$
et $V$ vérifiant les relations habituelles
\[
  \mathbf{F}\lambda = \lambda^{\sigma}\mathbf{F}, \qquad
  \lambda V = V\lambda^{\sigma}, \qquad
  \mathbf{F}V = V\mathbf{F} = p ,
\]
où $\lambda^\sigma$ est le Frobenius de $\lambda$ (la page écrit
$\lambda^{(p)}$ et $\lambda^p$). C'est le module de Dieudonné contravariant à
la Demazure--Gabriel.

\pagerange{60}{60}
\section*{Groupes finis sur une base de caractéristique $p$ (pages 60 à 65)}

\subsection*{Groupes finis étales (page 60)}

Soit $G$ un groupe fini étale annulé par $p^n$ sur $S$ de caractéristique $p$.
Pour tout épaississement $S'$ de $S$, $G$ s'étend de façon unique en un groupe
fini étale $G'$ sur $S'$, encore tué par $p^n$. Les
$G' \otimes_{\mathbf{Z}/p^n} \mathcal{O}_{S'}$ forment un cristal
$\mathbb{D}_*(G)$, et le Frobenius relatif $G \to G^{(p)}$, ici un
isomorphisme, donne un isomorphisme $F_* : \mathbb{D}_*(G) \to
\mathbb{D}_*(G)^{(p)}$. Par dualité, avec le dual étale
$\check G = \underline{\mathrm{Hom}}_{\mathrm{gr}}(G, (\mathbf{Q}_p/\mathbf{Z}_p)_S)$,
on pose $\mathbb{D}^*(G) = \mathbb{D}_*(\check G)$ : c'est un foncteur
contravariant, muni d'un isomorphisme $F^* : \mathbb{D}^*(G)^{(p)} \to
\mathbb{D}^*(G)$, et $V^*$ est déterminé par $F^*V^* = p$.

\emph{Cas où $G$ est plat sur $\Lambda_n = \mathbf{Z}/p^n$} (groupe de
Barsotti--Tate tronqué d'échelon $n$, étale). $\mathbb{D}^*(G)$ est un cristal
de $\Lambda_n$-modules localement libres de type fini, à $F$ isomorphisme (d'où
$V$), et l'on a une \emph{anti-équivalence} entre ces $G$ et ces
$F$-cristaux.\footnote{Une note en marge : « vérifier » (lecture douteuse). La
page écrit « équivalence » ; le foncteur étant contravariant, c'est une
anti-équivalence. Ce sont les $F$-cristaux « à racine unité » ; l'énoncé, pour
les systèmes locaux étales de $\mathbf{Z}/p^n$-modules, est la proposition 4.1.1
de N.~Katz, \emph{$p$-adic properties of modular schemes and modular forms}
(1973).} \emph{Cas général} : anti-équivalence entre groupes finis étales
$p$-primaires et $F$-$V$-cristaux localement (pour la topologie étale) sommes
de cristaux du type précédent. Il en résulte
$G(S) \simeq \mathrm{Hom}_{\mathrm{gr}}(\Lambda_{n,S}, G) \simeq
\mathrm{Hom}_{F\text{-}V\text{-cris}}(\mathcal{M}, \mathcal{O}_{S_n,\mathrm{cris}})$,
la formule s'interrompant au bord de la feuille.

\pagerange{61}{62}
\subsection*{Le triangle de dualité, cas étale et cas multiplicatif (pages 61 et 62)}

Les pages 61 et 62 poursuivent un argument commencé ailleurs : elles utilisent,
pour un groupe de Barsotti--Tate tronqué $G$, un « triangle exact de dualité »
qui relie le complexe co-Lie $\ell_G$ de $G$, la valeur $\Delta(G)$ du cristal
de Dieudonné sur $S$ (modulo $p$), et le dual décalé $\check\ell_{G^*}[1]$ du
complexe co-Lie du dual de Cartier $G^*$ :
\[
  \ell_G \longrightarrow \Delta(G) \longrightarrow \check\ell_{G^*}[1]
  \xrightarrow{\;+1\;} .
\]
C'est la forme dérivée de la filtration de Hodge
$0 \to \omega_G \to \mathbb{D}(G)_S \to \operatorname{Lie}(G^*) \to 0$.\footnote{Le
triangle n'est pas défini dans le dossier ; on le reconstitue d'après ses deux
cas particuliers. La théorie en est dans A.~Grothendieck, \emph{Groupes de
Barsotti--Tate et cristaux de Dieudonné} (Montréal, 1974), et L.~Illusie,
\emph{Déformations de groupes de Barsotti--Tate, d'après A.~Grothendieck},
Astérisque 127 (1985). La page le dessine en triangle, avec un signe
d'isomorphisme sur la flèche montante.}

\emph{$G$ étale.} Alors $\ell_G = 0$ et le triangle se réduit à un
isomorphisme $\Delta(G) \simeq \check\ell_{G^*}[1]$ : il ne donne aucune
structure en sus du cristal. On a
\[
  \check\ell_{G^*}[1] \simeq \Delta(G) \simeq \varepsilon^*(\mathcal{M}/p)
  = \mathbb{R}\underline{\mathrm{Hom}}_{\mathbf{Z}}(G, \mathcal{O}_S)
  \simeq \check G \overset{\mathbb{L}}{\otimes}_{\mathbf{Z}} \mathcal{O}_S
  \simeq (\check G \overset{\mathbb{L}}{\otimes}_{\mathbf{Z}} \mathbf{F}_p)
  \otimes_{\mathbf{F}_p} \mathcal{O}_S .
\]

\emph{$G$ de type multiplicatif} (annulé par une puissance de $p$). Il est
défini par son dual de Cartier $G^*$, étale, auquel s'applique ce qui
précède : $\mathbb{D}^*(G) = \Theta_*(G^*)$, c'est-à-dire
$\mathbb{D}^*(G)_{S'} = G^*_{S'} \otimes_{\mathbf{Z}_p} \mathcal{O}_{S'}$, où
$G^*_{S'}$ est le groupe étale sur $S'$ prolongeant $G^*$. Ici
$V : \mathbb{D}^*(G) \to \mathbb{D}^*(G^{(p)}) = \mathbb{D}^*(G)^{(p)}$ est un
isomorphisme, image par $\Theta_*$ du Frobenius de $G^*$, et $F$ est déterminé
par $V$. D'où des équivalences
\begin{quote}
groupes de Barsotti--Tate tronqués d'échelon $n$ de type multiplicatif
$\;\approx\;$ cristaux localement libres de type fini sur $S/\Lambda_n$ à $V$
isomorphisme (et $F$ nilpotent),
\end{quote}
et des équivalences analogues pour les groupes de type multiplicatif annulés
par une puissance de $p$.\footnote{« $\Lambda_n$ » et « nilpotent » sont de
lecture douteuse ; $F = pV^{-1}$ est bien nilpotent sur des cristaux de
$\Lambda_n$-modules.} Le triangle se réduit cette fois, puisque
$\check\ell_{G^*}[1] = 0$, à un isomorphisme $\ell_G \simeq \Delta(G)$, avec
$\ell_G \simeq G^* \overset{\mathbb{L}}{\otimes}_{\mathbf{Z}} \mathcal{O}_S
\simeq (G^* \overset{\mathbb{L}}{\otimes}_{\mathbf{Z}} \mathbf{F}_p)
\otimes_{\mathbf{F}_p} \mathcal{O}_S$.

\emph{Homomorphismes entre $p$-groupes étales et groupes de type
multiplicatif.} a) Un homomorphisme d'un groupe $G$ de type multiplicatif dans
un groupe étale $G'$ est nul ; et tout homomorphisme de $F$-$V$-cristaux (ou de
$F$-cristaux, ou de $V$-cristaux) de $\mathcal{M}' = \mathbb{D}^*(G')$ dans
$\mathcal{M} = \mathbb{D}^*(G)$ est nul : en écrivant la commutation à $F^n$,
où $F^n_{\mathcal{M}'}$ est un isomorphisme et $F^n_{\mathcal{M}}$ est nul pour
$n$ grand, on obtient le résultat pour les $F$-homomorphismes, et de même
pour les $V$-homomorphismes, les rôles étant échangés. b) Les homomorphismes
d'un groupe étale dans un groupe de type multiplicatif : le titre est posé, la
suite n'est pas écrite.\footnote{La page 63 ne porte que le transparent de la
page 64.}

\pagerange{64}{65}
\subsection*{Groupes étales et $F$-modules à racine unité (pages 64 et 65)}

Soit $S$ de caractéristique $p$, $\mathcal{C}$ la catégorie des schémas en
groupes finis étales $E$ sur $S$ annulés par $p$ — c'est-à-dire des faisceaux
étales de $\mathbf{F}_p$-espaces vectoriels localement libres de type fini —,
et $\mathcal{C}'$ celle des $\mathcal{O}_S$-modules localement libres $M$ munis
d'un isomorphisme $F_M : M^{(p)} \xrightarrow{\sim} M$. On a le foncteur
\[
  \varphi : \mathcal{C} \longrightarrow \mathcal{C}', \qquad
  E \longmapsto \bigl(E \otimes_{\mathbf{F}_p} \mathcal{O}_S,\;
  F_E \otimes_{\mathbf{F}_p} \mathcal{O}_S\bigr),
\]
le Frobenius relatif $F_{E/S}$ étant un isomorphisme de
$E \otimes \mathcal{O}_S$ sur $E^{(p)} \otimes \mathcal{O}_S \simeq
(E \otimes \mathcal{O}_S)^{(p)}$.

\textbf{Proposition.} \emph{Le foncteur $\varphi$ est une équivalence de
catégories.}

a) \emph{Pleine fidélité.} Par passage aux $\underline{\mathrm{Hom}}$, il
s'agit de voir que
$\Gamma(S, E) \xrightarrow{\sim} \{x \in \Gamma(S, E \otimes \mathcal{O}_S) \mid
F(x^{(p)}) = x\}$. Par descente on se ramène à $E$ constant, puis à
$E = \mathbf{F}_{p,S}$, et l'énoncé est l'exactitude de la suite
d'Artin--Schreier
\[
  0 \longrightarrow (\mathbf{Z}/p)_S \longrightarrow \mathbf{G}_{a,S}
  \xrightarrow{\;x \mapsto x^p - x\;} \mathbf{G}_{a,S} ,
\]
qu'il suffit de vérifier sur $\mathbf{F}_p$.

b) \emph{Essentielle surjectivité.} Soit $(M, F)$ ; il suffit de voir qu'il est
localement de la forme voulue. L'application
$\psi : M \to M$, $\psi(x) = F_M(x^{(p)}) - x$, est un épimorphisme étale de
schémas en groupes — sa différentielle est $-\mathrm{id}$ —, donc son noyau est
un schéma en groupes étale $E$ sur $S$. Pour qu'il soit fini, il suffit que son
$\mathbf{F}_p$-rang soit localement constant, égal au rang de $M$ ; et pour que
$E \otimes \mathcal{O}_S \to M$ soit un isomorphisme, il suffit de le voir fibre
par fibre géométrique. On est ramené au

\textbf{Lemme.} \emph{Soit $M$ un espace vectoriel de dimension finie sur un
corps $k$ séparablement clos de caractéristique $p$, muni d'un isomorphisme
$F_M : M^{(p)} \xrightarrow{\sim} M$, et
$E = \{x \in M \mid F(x^{(p)}) = x\}$. Alors
$E \otimes_{\mathbf{F}_p} k \xrightarrow{\sim} M$.}\footnote{La page écrit
d'abord « alg.\ clos », corrigé de sa main en « sép.\ clos », qui suffit : les
équations de $E$ sont séparables. Le lemme n'est pas démontré dans le
dossier ; c'est un cas du théorème de Lang (1956), et l'ensemble de la
proposition est la proposition 4.1.1 de Katz (1973) déjà citée, au niveau
$n = 1$.}

\emph{La version de Witt.} Les notes amorcent ensuite la même équivalence au
niveau $p^n$, avec les vecteurs de Witt : $\mathcal{C}$, les schémas en
$p$-groupes finis étales sur $S$ (faisceaux localement constants de $p$-groupes
finis) ; $\mathcal{C}'$, les faisceaux de $W_S$-modules (pour la topologie de
Zariski — « Zar ou fppf, c'est kif-kif », dit la marge) localement isomorphes à
une somme de $W_{n,S}$, munis d'un isomorphisme $F_M : M^{(p)} \xrightarrow{\sim}
M$ ; $\mathcal{C}''$, les schémas en $W_S$-modules lisses. Les foncteurs sont
$\varphi(E) = E \otimes_{\mathbf{Z}_p} W_S$ avec le $F$ évident, vers
$\mathcal{C}'$ ; $\psi(E) = E \otimes_{\mathbf{Z}_p} W_S$, vers $\mathcal{C}''$ ;
et une équivalence $i : \mathcal{C}' \to \mathcal{C}''$, qui associe à $M$ le
schéma représentant le faisceau\ldots\ La note s'arrête là, le reste de la page
est vide ; le triangle $\psi = i \circ \varphi$ qu'elle dessine est
annoncé, non démontré.\footnote{L'exposant de $W_S$ dans la définition de
$\mathcal{C}''$, lu « fl » (topologie plate), est douteux. La page 66 est une
feuille blanche à en-tête.}

\end{document}
